% Declared-model duty 9. Independent author derivation saved before E4 intake.
% CURRENT BINDING SCOPE -- full GOAL_alt.md read on 2026-10-15.
% The binding model uses continued occupied spectral P[h], [P,h]=0,
% the compatible -J row, explicit spherical source projection, and
% the full source adjoint in the site Hamiltonian. Its reduced quantum
% prescription is expansion about a positive Hartree mass with full
% Wick ordering against that same continued reference. The pending
% BG7/QR22--24 POINT-gradient compact promotion is withdrawn: those
% are control theorems and their mismatch is absent for this coupling.
% IR/BG transported-representation/variance claims retain the literal
% transported-reference domains actually proved, not the new spectral
% reference. The strong full-Fock Gauss obstruction stands as delivered;
% no further continuation of it is undertaken. Existing symmetric
% source Hessians, spectral contact algebra and full-Wick identities
% are candidates for the revised reduced construction at their stated
% branch/domain premises. The compact now gives the constructive
% RR reduction; the earlier control proofs remain in the comments.
% STATUS [R]: constructive continued-spectral/adjoint Hartree derivative,
% closed radial/native response and full-reference-Wick quadratic
% correction. The complete ordered positive-mass domain remains open.
\noindent\textbf{[R] Reduced spherical boundary Hamiltonian.}
With $G$ defined by the Cavendish measurement, use lattice units,
$\mu=ma$ and $g=G/a^2$.
Fix a finite native parent, its energy/current vertices $B_I$ and actual
spherical averaging map $W$. The site generator is
$h=\sum_I[W^\dagger(N,v)]_I B_I$.
On a regular positive-mass, positive-clock branch with coherent $-J$,
$u=v/N<0$, and a differentiable continued occupied projector $P$,
\[
 (N\delta M)'=4\pi r^2N(\delta\rho+u\delta J),\qquad
 \delta E=\operatorname{Tr}[h\,\delta(C-P)]
          =\operatorname{Tr}(h\,\delta C).
\]
The last equality follows from $[h,P]=0$ and projector tangency.
Thus $i\dot C=[h[C],C]$ conserves the closed-system boundary charge;
$\dot D=-i[h,D]-\dot P$ for $D=C-P$, so nonadiabatic pairs source it.
The local chart needs separated opposite occupations and a regular
implicit inverse; zero-energy crossings alone do not change occupations.

The closed radial density--current Hessian $K$ is given below, with
$\delta^2E=\int dr[8\pi grN\,\delta M\delta J/u
-4\pi g^2NJ(\delta M)^2/u^3]$.
Put $A=W^\dagger K W$ and
$\Pi_{IJ}=\operatorname{Tr}(B_I\,DP_h[B_J])$.
The symmetric native covariance Hessian is
\[
 K_c=(I+A\Pi)^{-1}A,\qquad
 K_c=\sqrt g\,A_{1/2}
       +g(A_1-A_{1/2}\Pi_0A_{1/2})+O(g^{3/2}).
\]
The second formula retains the order-$g$ reference contact on a regular
weak-field branch. Define
$\mathcal L(X)=\sum_{IJ}(K_c)_{IJ}B_I\operatorname{Tr}(B_JX)$.
At a Hartree background $\bar C$, $\bar P=P(\bar h)$,
$\bar D=\bar C-\bar P$, the fully reference-Wick-ordered expansion is
\[
 \begin{aligned}
 \widehat H^{(2)}
 &=c_H+{:}d\Gamma(\bar h-\mathcal L(\bar D)){:}_{\bar P}
 +\frac12\sum_{IJ}(K_c)_{IJ}
       {:}d\Gamma(B_I)d\Gamma(B_J){:}_{\bar P},\\
 c_H&=\bar E-\operatorname{Tr}(\bar h\bar D)
           +\tfrac12\operatorname{Tr}[\bar D\mathcal L(\bar D)].
 \end{aligned}
\]
Its Hartree expectation is the quadratic Taylor expansion of $E[C]$;
the derivative at $\bar C$ is $\bar h$. For Gaussian $C=\bar P+D$,
\[
 E_F=-\frac12\sum_{IJ}(K_c)_{IJ}\operatorname{Tr}(B_IDB_JD),
 \qquad h_F=-\sum_{IJ}(K_c)_{IJ}B_IDB_J .
\]
All reference contractions are removed. Pair energy and its derivative
vanish at $D=0$; for rank-one particle $D$, Hartree self-attraction cancels
Fock exactly. In the slow-particle, small-$\ell/L$ limit the spherical
kernel is $-g\mu^2/\max(r,s)$; the displayed response retains finite
averaging and current terms.

For the full mass-response vertex $A_r^{\rm eff}=\nabla_C M_r$ below,
$V_r=\operatorname{Tr}[\bar C A_r^{\rm eff}(1-\bar C)A_r^{\rm eff}]$.
Report its sea value and signed excess separately.
The positive rest-mass/Slater number approximation gives relative
particle fluctuations $\leq N_r^{-1/2}$, with $N_r$ the expected
enclosed particle number; a positive bounded rest-mass kernel gives
$|E_F|/|E_H|\leq k_{\max}/(k_{\min}N_f)$.
Neither replaces the full variance or general current kernel.
For task 8, after the one-body Wick correction a centered-density
implementation still needs
$k(P)=-\operatorname{Tr}(hP)-\frac12\sum_{IJ}V_{IJ}
\operatorname{Tr}(PB_IPB_J)$, hence $k(P_N)-k(P_0)$ in sector energies.
This scalar changes no fixed-sector eigenvector or entropy; clipped
sea-subtracted variance is not physical variance.

Every finite-order operator is self-adjoint on a finite parent;
its controlled use is about the positive Hartree branch.
The full ordered series' invariant domain, parent limit and uniform
allowed-$\ell$ control remain open.
Angular/radial quantum fluctuations exclude a strong operator Gauss law
on the full Fock space; the construction here is the reduced one.

% INDEPENDENT DERIVATION / CHECK DETAIL.
% No icloud E4 or new solve exchange numbers were read before this file
% was first saved. Same-lane RED was assigned its own independent route.
% Geometry is hydrodynamic and has no lattice-scale components. G is
% defined by the Cavendish measurement. The continuum Gauss law is
% declared; no microscopic geometric stencil or metric phase space is used.
% Lattice-unit reporting uses g=G/a^2, mu=ma, R=r/a, L=ell/a,
% htilde_0=a h_0, rhotilde_x=a rho_x and ktilde=a k. Thus
% a E_H=-(g/2)sum ktilde rhotilde_x rhotilde_y,
% a E_F=(g/2)sum ktilde Tr(etilde_x D etilde_y D), and
% a^2 Var M=Tr[C Atilde(1-C)Atilde]. The radial identity is
% u^2=2g Mtilde/R with Mtilde=a M. These are functions of mu,g,
% the specified state and the explicitly tested L; no second coupling
% or geometric lattice parameter is introduced.
%
% 1. Newton reduction and the half factor.
% In the weak, slow static sector the continuum Einstein Gauss law fixes
% the ordinary Newton potential. In flat slicing v=O(sqrt(G)), and the
% local normal-frame kinetic energy contains the compensating drift.
% At leading order a static particle has normal momentum p=-m v.
% Completing p^2/(2m)+v.p gives (p+m v)^2/(2m)-m v^2/2.
% The irrotational phase transformation removes p+m v at this order.
% The remaining scalar potential is m(delta N-v^2/2). The interior
% constraint solution has v^2=2GM(r)/r and
% delta N=-G int_r^infinity dM(s)/s, yielding the displayed Phi.
% Equivalently solve the continuum Poisson law with zero potential at
% infinity. Spherical integration of 1/|x-y| gives 1/max(r,s).
% The boundary binding energy is
% U=-G int_0^infinity M(r)dM(r)/r
%  =-(G/2) int dr ds M'(r)M'(s)/max(r,s)
%  =-(G/2) int_0^infinity M(r)^2 dr/r^2.
% It has functional derivative Phi. The energy sum int rho Phi is 2U.
% This also explains why mere substitution of density-dependent N,v into
% H_m is not a derivation of the reduced many-body Hamiltonian: its
% rearrangement term and the boundary gravitational energy matter.
% This reduction requires the compatible momentum sign -J; it is not a
% consequence of the literal +J row. Independently postulating the Poisson
% potential would give the same Newton functional, but would not repair
% that row. This is a weak slow-matter result. At relativistic
% velocities, momentum, stress and retardation terms cannot be replaced
% by rho rho without an actual reduction of the declared continuum law.
%
% 2. Coarse graining and a closed lattice kernel.
% W_ell(x)=(2pi ell^2)^(-3/2) exp[-|x|^2/(2ell^2)]. Source averaging is
% convolution with W_ell followed by angular projection. On the continuum
% parent these commute with the Newton Green operator. Hence the symmetric
% one-smearing pair kernel is angular-average[erf(d/(sqrt2 ell))/d],
% d=|x-y|. Substitution d dd=-rs dcos(theta) gives the displayed F_ell.
% At r=0, k_ell(0,s)=erf(s/(sqrt2 ell))/s; at r=s=0 it is
% sqrt(2/pi)/ell. At ell->0 and nonzero radii it tends to 1/max(r,s).
% The kernel is evaluated at actual lattice sites r_x=|a x|; it is not
% a shell entropy surrogate. A smooth source with width much larger than
% a can be represented by these site charges and its resolved continuum
% convolution, with sampling errors checked separately.
%
% There is an important averaging count. Squaring the physical averaged
% mass M_ell in int M_ell^2/r^2 would smear BOTH charges and yield
% erf(d/(2ell))/d. Its Hartree derivative has TWO source filters and so
% does not reproduce a once-averaged source in the pointwise matter
% Hamiltonian. The chosen kernel has exactly one source filter. It can
% also be factored as int F_{ell/sqrt2,x}(r) F_{ell/sqrt2,y}(r)dr/r^2,
% where F_{b,x}(r)=int_{|z|<r}W_b(z-a x)d^3z. This factorization proves
% positive definiteness and supplies a symmetric mass-square ordering;
% its half-filter mass is a mathematical factor, not the physical source
% mass M_ell. If both matter coupling and sources are instead coarse
% grained, the two-filter convention is a different explicit prescription.
% They agree only to the controlled long-wavelength accuracy.
%
% 3. Ordering and Wick derivation.
% P_0=1_{h_0<0} is the reference at G=0 for this first-order pair term.
% At nonzero fields the declared reference is the transported incoming sea
% P_T, not instantaneous P_-(h). For a fixed finite-duration smooth weak
% Hartree history, Duhamel expansion gives P_T-P_0=O(sqrt(G)) in native
% PG scaling when h-h_0=O(sqrt(G)). Inserting that difference into an
% explicit O(G) quartic term is higher order at fixed duration. Uniformity
% at late times is not inferred. Other reference contacts must still be
% computed. The spectral susceptibility in sections 8--9 is a separately
% labeled instantaneous-ground comparison, not the transported law.
% rho_x in this file is energy per cell, not energy density; divide by
% a^3 to compare to the continuum Gauss-law source convention. This is
% why no extra a^6 appears in the two sums in the displayed pair energy.
% Define q_x=c^dag e_x c-tr(P_0 e_x). Declared full reference Wick
% ordering of the four fermion operators, followed by Hermitian
% symmetrization x<->y, removes reference contractions and same-particle
% self-interaction. It is a separate ordering choice: multiplying only
% the already mean-subtracted q_x operators does NOT perform it.
% For a gauge-invariant Gaussian covariance C, Wick's theorem then gives
% < :q_x q_y:_{P_0} >=tr(e_x D)tr(e_y D)-tr(e_x D e_y D), D=C-P_0.
% With -(G/2)sum k_xy this proves E_H,E_F. Symmetry of k_xy makes the
% apparently unsymmetrized trace and h_F real/Hermitian after summation.
% Differentiation gives h_H=-G sum_xy k_xy rho_y e_x and
% h_F=+G sum_xy k_xy e_x D e_y. These are finite lattice matrices with
% the original four orbitals; no additional flavor count is implicit.
% The exact rest-energy sector uses positive particles only and
% e_x=m Pi_x+O(p^2/m). Then the Fock kernel is the displayed Hadamard
% product. If both particles and holes occur, or energy-density vertices
% are used outside this sector, D need not be positive and a universal
% repulsive sign or N_f power cannot be inferred from particle counting.
%
% 4. Relative size and the actual spherical restriction.
% In the positive-particle Slater sector D=Gamma, Gamma^2=Gamma,
% tr Gamma=N_f. With k_min<=k_xy<=k_max on the occupied support,
% 0<=E_F<=(Gm^2/2)k_max tr Gamma^2,
% |E_H|>=(Gm^2/2)k_min N_f^2, so
% E_F/|E_H| <= (k_max/k_min)/N_f. This bound explicitly retains the
% profile/central-resolution dependence; k_max is O(1/ell), so it is
% not a uniform shape-independent constant as ell/R->0.
% The matching lower bound is k_min/(k_max N_f), because the same
% nonnegative exchange weights sum to tr Gamma^2=N_f. In the unsmoothed
% annulus R_min<=r<=R_max these become the exact two-sided bounds
% R_min/(R_max N_f)<=E_F/|E_H|<=R_max/(R_min N_f), with no assumption
% on exchange tails. The independent RED established the matching lower
% bound; the author checked it from the same positive matrix entries.
% For a smooth macroscopic profile whose exchange weight is dominated by
% separations small compared with R and whose tail is controlled, replace
% k_xy by k_0(r_x,r_x)=1/r_x over the exchange correlation range.
% Sum_y tr Gamma_xy Gamma_yx=tr Gamma_xx then gives
% E_F=(Gm^2/2) int n(x)/r d^3x plus correlation/smoothing corrections.
% A uniform ball of radius R has E_H=-3Gm^2 N_f^2/(5R) and
% E_F=3Gm^2 N_f/(4R), hence ratio 5/(4N_f). Boundary, central and
% smoothing corrections vanish in the stated resolved large-object
% limit; this is not an arbitrary-state equality.
% A zero-temperature Fermi projector has algebraic rather than exponential
% decay. The approximation is an approximate-identity limit for |Gamma|^2,
% not a finite correlation-length assertion. For the bulk continuum Fermi
% gas its normalized tail outside a separation d is O(1/(p_F d)); a
% first-moment estimate can carry log(p_F R)/(p_F R). Such tails, the
% boundary region and the integrable r=0 weight must be controlled in
% using the displayed uniform-ball coefficient.
% If one instead retains the full Newton kernel, the usual homogeneous
% Fermi-gas exchange integral is proportional to Gm^2 int n^(4/3), so
% at fixed shape E_F/|E_H| is proportional to N_f^(-2/3), with its spin
% degeneracy factor. This unsmoothed result also describes a filtered full
% kernel only if ell p_F<<1. For ell p_F>>1 its smooth diagonal instead
% gives E_F of order Gm^2 N_f/ell; thus no unsmoothed Fock scaling is
% silently claimed independent of the source resolution. Identical
% spherical Hartree potentials do not make
% these quantum theories equivalent. No such angular geometric modes
% are silently added to the present spherical-only model.
%
% 5. Exact fluctuations and the negative-spectrum obstruction.
% F_r is multiplication by the cumulative Gaussian ball weight
% F_r(x)=int_{|z|<r} W_ell(z-a x)d^3z, and A_r={F_r,h_0}/2.
% M_r=c^dag A_r c-tr(P_0 A_r). For every Gaussian C,
% <M_r>=tr[(C-P_0)A_r],
% Cov(M_r,M_s)=Re tr[C A_r(1-C)A_s] (symmetrized covariance).
% Mean subtraction changes neither variance nor this covariance.
% In the positive nonrelativistic particle approximation A_r=m F_r,
% Var=m^2[tr(F_r^2 Gamma)-tr(F_r Gamma F_r Gamma)]
%     <=m^2 tr(F_r Gamma)=m^2 N_r, since 0<=F_r<=1.
% For a sharp resolved count, <M>=m N_r, giving the stated relative
% bound. For a smooth count the same mean definition works and the bound
% remains valid. Fixed-total-number Slater states have zero fluctuation
% of the full particle mass when F_r=1, rather than exactly sqrt(N_f).
% Number alone gives no bound for arbitrary states: a fixed-N_f coherent
% superposition of orthogonal states with all particles well inside or
% well outside the counting surface has <M_r>=mN_f/2 and
% Var M_r=m^2 N_f^2/4 in the resolved rest-mass limit. Its relative
% fluctuation is 1, not N_f^(-1/2). The localized orbitals can be chosen
% smooth and much wider than a and 1/m, and far from the averaging collar.
% Thus the Slater/controlled-correlation premise is a physical state
% restriction, not a consequence of the large particle count by itself.
%
% On the actual four-component sea the vacuum term is generally nonzero:
% Var_0 M_r=tr[P_0 A_r(1-P_0)A_r]>=0. In orthonormal momentum basis,
% A_r(p,q)=F_r(p,q)[h_0(p)+h_0(q)]/2. Therefore its exact contribution is
% (1/4)sum_{pq}|F_r(p,q)|^2(E_p-E_q)^2
%       *[1-(m^2+d_p.d_q)/(E_p E_q)], d_i=sin(ap_i)/a.
% On a massless finite parent this divided formula applies to E_p E_q>0;
% zero-mode contributions use the actual chosen filling separately.
% Generic nonconstant F_r connects unequal energies and unequal normalized
% Dirac symbols, making the sum strictly positive. This is an explicit
% lattice counterexample with N_f=0 to a universal N_f^(-1/2) assertion.
% A positive semidefinite M_r with <0|M_r|0>=0 would satisfy M_r|0>=0
% and hence Var_0=0. The positive variance proves a negative spectral
% subspace of the mean-subtracted self-adjoint M_r. Thus a self-adjoint
% v satisfying v^2=(2G/r)M_r cannot exist on the full matter Fock space
% for the direct leading spherical mass identification. An ordering of
% the square root cannot turn a negative operator into a positive one.
%
% A constructive allowed Hartree expansion is available at a state with
% bar M_r>0 and controlled delta M_r/bar M_r:
% v_r=-sqrt(2G bar M_r/r)[1+delta M_r/(2bar M_r)
%                         -delta M_r^2/(8bar M_r^2)+...].
% Powers of the same M_r are self-adjoint; mixed radial/current products
% can be Weyl ordered explicitly. The state-domain test is the actual
% covariance above, including its reference term. To leading order
% Delta v/|bar v|=Delta M/(2bar M); the mean shift is
% -Var(M)/(8bar M^2) relative to bar v. Chebyshev bounds the weight with
% M_r<=0 by Var(M_r)/bar M_r^2. Convergence in operator norm requires
% a stronger spectral support condition; a variance bound alone gives
% only a state-specific perturbative control, not a full-Hilbert theorem.
%
% 6. Averaging-scale test and remaining requested scope.
% In the once-filtered Hartree functional for a smooth spherical density,
% U_ell=-(G/2)int d^3k/(2pi)^3 (4pi/k^2)
%                         exp(-ell^2 k^2/2)|rho(k)|^2.
% Consequently U_ell-U_0=pi G ell^2 int rho^2 d^3x+O(ell^4 grad rho^2).
% Its leading fractional correction is O(ell^2/R^2), not identically zero.
% On the diagonal k_ell(r,r)=1/r-ell/(sqrt(2pi)r^2)
% up to exponentially small terms for r>>ell. Exchange may therefore
% retain an O(ell/R) relative correction to its already O(1/N_f) term.
% The vacuum fluctuation is an explicitly ell-dependent observable and
% requires its own comparison to bar M^2; mean subtraction cannot erase it.
% Vary ell in the allowed hierarchy and report convergence/tolerances.
% Exact ell independence at finite resolution is not a consequence of the
% declarations. The source/matter filtering convention must be specified
% before comparing a numerical coefficient.
%
% The complete relativistic pair operator is not fixed here by a Newton
% ansatz. Constraint equations specify sources and instantaneous fields;
% their operator substitution alone does not specify a reduced symplectic
% form, a boundary Hamiltonian, all stress terms, or a positive-mass domain.
% Required ways forward are: derive that reduction in the stated spherical
% theory; use the explicitly controlled weak slow-particle sector above;
% or compute the actual state-specific Hartree expansion with its Fock
% and full mass variance, retaining any vacuum contribution. This preserves
% the full duty rather than silently identifying it with a rest-mass model.
%
% 7. General source reduction, before E4 intake.
% This continuation was derived independently after the weak-sector note,
% using the actual declared_h2_continuum.tex integrating-factor construction.
% It does not silently change the literal momentum row. Introduce
% sigma=+1 for the compatible momentum RHS -J and sigma=-1 for literal +J.
% Let dV_r=4pi r^2 dr, u=v/N have a fixed nonzero sign for r>0,
% M=ru^2/(2G)>0, M(0)=0, N(infinity)=1, E=M(infinity). The radial equations
% are M'=4pi r^2(rho+sigma u J), (log N)'=-sigma 4pi Gr J/u.
% The linearized identity is
% (N delta M)'=4pi r^2 N(delta rho+sigma u delta J).
% Hence delta E=int dV (N delta rho+sigma v delta J). With literal +J,
% this does not give the declared +v.P matter Hamiltonian. Already at J=0,
% delta N(r)/delta J(s)=-G Theta(s-r)/(s u(s)) whereas
% delta v(s)/delta rho(r)=+G Theta(s-r)/(s u(s)), using dV source measure.
% Thus independent-source Maxwell reciprocity fails for the proposed (N,v)
% gradient with that literal sign. This is a source-functional obstruction;
% realizability of arbitrary source tangents by an implicit covariance
% branch is a separate question and is not presumed. The compatible sign
% restores this necessary reciprocity test. Every formula below explicitly
% adopts sigma=+1 as an allowed comparison, not an authorized model repair.
%
% Fix a smooth compatible-sign background with positive M and N and with
% all following integrals finite. Set n=log N. Along a linear source path,
% delta u=G delta M/(r u),
% delta n'=-4pi Gr delta J/u+4pi G^2 J delta M/u^3.
% Differentiating the boundary variation and integrating by parts gives
% delta^2 E=int dr [8pi Gr N delta M delta J/u
%                         -4pi G^2 N J (delta M)^2/u^3].
% The endpoint term [delta n N delta M]_0^infinity vanishes by the stated
% center and asymptotic conditions. This exact source Hessian includes the
% full normal energy and radial current, without a rest-mass approximation.
% At J=0 its pure-rho block vanishes: the PG normal-source boundary charge
% is E=int rho dV on that branch. Newton binding therefore cannot be read
% from this block without the static-particle/canonical-variable conversion
% used in section 1. Marginally bound comoving matter is a different state.
%
% Closed response blocks in the dV measure follow from
% delta M(r)=int dV_s Theta(r-s) [N_s/N_r]
%                             [delta rho_s+u_s delta J_s].
% Define I(R)=int_R^infinity 4pi G^2 J_r/(N_r u_r^3) dr and R=max(s,t).
% Then, with q=(rho,J), delta phi_alpha(s)=int dV_t K_alpha beta(s,t)
% delta q_beta(t), phi=(N,v), and
% K_rho rho(s,t)=-N_s N_t I(R),
% K_rho J(s,t)=G Theta(t-s) N_s/(t u_t)-u_t N_s N_t I(R),
% K_J rho(s,t)=K_rho J(t,s),
% K_J J(s,t)=G Theta(s-t) N_t u_t/(s u_s)
%           +G Theta(t-s) N_s u_s/(t u_t)-u_s u_t N_s N_t I(R).
% Substitution in int dV_s dV_t delta q(s) K(s,t) delta q(t) reproduces
% the Hessian above. Theta(0)=1/2 is a symmetric continuum convention;
% actual finite cells must integrate the kernel with their source weights.
% This is a closed state-dependent radial kernel, not an added angular
% gravitational mode or a postulated lattice-scale geometric law.
%
% The expansion at fixed normal sources is not an ordinary density-pair
% expansion in G. For G-independent smooth rho,J and M_0(r)>0, put
% z_0(r)=sign(u) sqrt(2M_0(r)/r), M_0'=4pi r^2 rho.
% On the regular branch M=M_0+sqrt(G)M_1+G M_2+O(G^(3/2)),
% M_1(r)=int_0^r dV_s z_0(s)J_s,
% M_2'=4pi r^2 J_r M_1(r)/(r z_0(r)). Therefore
% E=E_0+sqrt(G) int dV_r z_0(r)J_r
%   +(G/2) int dV_r dV_s K_JJ^(1)(r,s) J_r J_s+O(G^(3/2)),
% K_JJ^(1)(r,s)=z_0(min(r,s))/[max(r,s) z_0(max(r,s))]
%             =sqrt[M_0(min)/M_0(max)]/sqrt(rs).
% This expression is for nonzero r,s, extended only where the regular-center
% integrals converge. Its coefficient depends on rho through M_0. It also
% has the preceding nonpolynomial sqrt(G) term. The literal-sign branch
% reverses the sqrt(G) term but leaves this quadratic current term unchanged.
% In a static weak object J=O(sqrt(G)); that is a different scaling and
% recovers the Newton sector after transforming the normal kinetic energy.
% No universal, background-independent relativistic O(G) rho-rho operator
% follows by retaining only the displayed current pair and dropping the
% density dependence or the lower-order flow term.
%
% 8. Instantaneous spectral-reference comparison and local quantum ordering.
% This section assumes P=P_-(h), explicitly a spectral comparison rather
% than the binding transported reference of section 14. Assume an actual
% smooth solution branch of that implicit Gauss+reference
% system, not arbitrary source functions declared realizable by matter.
% At finite lattice parent, let B_I be the original site energy and radial
% current vertices, including the continuum radial source-assignment weights.
% In finite source coordinates write h=sum_I phi_I B_I,
% q_I=Tr[B_I(C-P(h))], P=1_{h<0}. Volume measures are absorbed consistently
% into these coordinates; there is no extra flavor or cell-volume factor.
% Require a fixed-rank nonzero spectral gap and a differentiable branch.
% Tr(h delta P)=0 then gives delta E=Tr(h delta C) in the compatible,
% unsmoothed construction. All derivatives of P must nevertheless remain
% in the second variation. Its Frechet derivative defines
% Pi_IJ=Tr[B_I DP_h[B_J]]
% =2 Re sum_{n occupied,m empty}
%    <n|B_I|m><m|B_J|n>/(epsilon_n-epsilon_m).
% This matrix is real symmetric and negative semidefinite, as follows by
% contracting it with a real vertex combination B=sum b_I B_I.
% It is an exact finite-matrix formula for the declared sine operator and
% its two-step current, not a continuum dispersion replacement.
% Put b_I(delta C)=Tr(B_I delta C). Linearization gives
% delta q=b(delta C)-Pi delta phi, delta phi=K delta q.
% If I+K Pi is invertible on the allowed source tangent space,
% delta phi=K_eff b(delta C),
% K_eff=(I+K Pi)^(-1)K=K(I+Pi K)^(-1).
% K_eff is symmetric whenever this inverse exists. A zero eigenvalue or a
% closing reference gap invalidates the local reduction; neither is assumed
% absent merely because the Hartree mass is positive.
% Equivalently, differentiating Tr(h delta P)=0 gives
% Tr(h delta^2 P)=-Tr(delta h delta P), so the pulled-back covariance
% Hessian equals the source Hessian plus Tr(delta h delta P), and its
% gradient response is delta h=sum_I B_I (K_eff b)_I.
% This is the Hessian in affine covariance coordinates. Along a curved
% pure-Slater path the ordinary second derivative also contains the
% kinematical term Tr(h delta^2 C). A particle-hole tangent restriction must
% be applied to these expressions, not replaced by arbitrary independent
% rho,J variations. The susceptibility describes reference feedback; it
% does not renormalize G or screen the fixed exterior Gauss-law charge.
% In particular, if K=sqrt(G) K_(1/2)+G K_1+O(G^(3/2)) and Pi=Pi_0+O(sqrt(G)),
% then K_eff=sqrt(G) K_(1/2)
%                +G[K_1-K_(1/2) Pi_0 K_(1/2)]+O(G^(3/2)).
% The reference contact is thus already an O(G) part of the response in
% these variables. Keeping the bare current kernel alone misses it.
%
% A concrete local fluctuation quantization can now be declared at a fixed
% accepted background (C_bar,P_bar,K_eff), if that branch exists. Freeze
% these coefficients, use qhat_I=c^dag B_I c-Tr(P_bar B_I), and fully
% reference-Wick order (1/2)sum K_eff,IJ qhat_I qhat_J; subtract the linear
% and constant Hartree terms when centering about qbar. For a Gaussian C,
% with D=C-P_bar, the uncentered quadratic expectation is
% E_H^(2)=(1/2)sum K_eff,IJ Tr(B_I D)Tr(B_J D),
% E_F^(2)=-(1/2)sum K_eff,IJ Tr(B_I D B_J D),
% h_F^(2)=-sum K_eff,IJ B_I D B_J.
% The same formulas with the appropriate linear/constant centering have
% the identical exchange term. They exhibit the exact finite-lattice
% density-current Fock kernel for this specified quadratic expansion.
% Its sign and size are state dependent because K_eff and D need not be
% positive. They do not imply the slow-particle 1/N_f estimate for arbitrary
% relativistic states. The exact variance remains the covariance formula
% in section 5 with the actual energy/current vertices and actual P_bar.
% Varying the background varies K_eff and generates higher/rearrangement
% terms. Thus this is a state-specific quadratic boundary-energy expansion,
% not a globally defined state-independent operator Hamiltonian or a
% proof of the literal binding model. Ordering alone supplies neither the
% positive-mass domain nor the missing reduced construction.
%
% The corresponding general mass fluctuation is also constructive at this
% linearized level. Write L_r for the source-to-mass row in section 7:
% delta M_r=L_r delta q. Its covariance pullback is
% delta M_r=L_r(I+Pi K)^(-1)b(delta C).
% Define the Hermitian native vertex
% A_r^eff=sum_I [L_r(I+Pi K)^(-1)]_I B_I.
% The linearized centered operator is c^dag A_r^eff c-Tr(C_bar A_r^eff),
% so its exact Gaussian variance is Tr[C_bar A_r^eff(1-C_bar)A_r^eff].
% This includes the radial current and the actual reference response, unlike
% the direct energy-only mass in the domain counterexample of section 5.
% It is the variance of the declared linearized operator; nonlinear mass
% terms add higher corrections whose ordering and control must be checked.
% With source averaging and native field readout, let W map native charges
% to smooth continuum sources, R map continuum fields to the native point
% coupling, and set A=R K W. Replace the mass row by
% L_r W(I+Pi A)^(-1). This linearized constraint-response formula still
% requires a regular branch; it does not certify Hamiltonian integrability.
% The corresponding field variances use vertices sum_J T_IJ B_J with
% T=(I+A Pi)^(-1)A. Actual ratios to positive Hartree M,N,v, rather than
% particle number by itself, control this fluctuation expansion.
%
% 9. Finite averaging and the exact admissible-variation criterion.
% For qbar=W_ell q in continuum source coordinates, the compatible boundary
% functional is E[W_ell q], whose bare-source gradient is W_ell^dagger phi;
% it is not generally phi. Native source/readout maps are restored below.
% In the instantaneous-reference comparison, write
% hbar=h[W_ell^dagger N,W_ell^dagger v] and P=P(h). Its variation is
% delta E=Tr(hbar delta C)-Tr[(hbar-h)delta P].
% Equality to the declared h requires
% Tr[(hbar-h)delta C]=Tr[(hbar-h)delta P]
% for EVERY admissible tangent. W_ell^dagger phi=phi is sufficient, not
% necessary: a separately established reference-contact compensation on
% the actual solution branch has not been excluded. The displayed mismatch
% alone is therefore not a no-go theorem on all covariance restrictions.
%
% There is, however, an independent-source reciprocity obstruction to
% simply replacing the boundary gradient by the point-field phi. At J=0,
% N=1, define A=K_rho J by (Af)(r)=int_r^infinity c(s)f(s)ds,
% c(s)=4pi Gs/u(s). Gaussian scalar and radial-vector averages are
% W_0=exp[(ell^2/2)L_0], W_1=exp[(ell^2/2)L_1],
% L_0=partial_r^2+(2/r)partial_r, L_1=L_0-2/r^2.
% Reciprocity of the point-field response K W would require W_0 A=A W_1.
% But partial_r(L_0 A-A L_1)=-[L_1,c]
% =-2c' partial_r-c''-2c'/r, generally nonzero.
% For example a smooth exterior interval has u proportional to r^(-1/2),
% hence c proportional to r^(3/2) and c' nonzero. Compactly supported test
% sources inside that interval establish nonzero antisymmetric response at
% O(ell^2), within a<<ell<<r. This excludes an independent-source potential
% for that point-coupling prescription at finite ell. It does not prove
% that every restricted implicit covariance branch violates the exact
% criterion above. An adjoint-filtered matter generator with its reference
% rebuilt for that same generator is a constructive alternative model;
% the unsmoothed compatible branch is another controlled comparison.
% Neither is silently substituted for the binding point-field prescription.
%
% The finite-average limitation is visible even in the slow static sector.
% Let rho_0=m times particle number per normal-slice volume, held fixed;
% it is not the comoving proper rest density. At leading weak static order,
% rho=rho_0(1+u^2/2)+O(G^2), J=-rho_0 u+O(G^(3/2)),
% u^2=2G M_{0,ell}(r)/r, with M_{0,ell} the integral of W rho_0.
% Integrating the compatible averaged radial row M'=4pi r^2(W rho+u WJ)
% and using normalized W and its physical-volume adjoint gives
% E-E_0=(1/2)int dV rho_0 u^2-int dV rho_0 u (W^dagger u)+O(G^2).
% The scalar and radial-vector averages must retain their actual placements.
% At W=1 this becomes -G int dV rho_0 M_0/r, the symmetric Newton energy.
% At finite ell, even -G int dV rho_0 M_{0,ell}/r alone is not generally
% the once-filter symmetric pair: its inner/outer shell terms need not
% agree after averaging just one source profile. The displayed extra
% u(W^dagger u-u) term is also generally nonzero. For a fixed smooth
% profile they are controlled gradient corrections as ell/R tends to zero,
% not exactly absent at every allowed ell. Accordingly sections 1--4 give
% an explicit once-filter Newton prescription with the correct Newton
% Hartree derivative and exchange; they do not prove its equality to the
% finite-ell boundary reduction of the unaltered point-field model.
% This direct static substitution was checked independently by the RED.
% A smooth compact rho_0>=0 with regular center makes the discrepancy
% explicit. For the Gaussian coordinate variance ell^2, write
% u_0^2=2GM_0/r and M_{0,ell}-M_0=2pi ell^2 r^2 rho_0'+O(ell^4).
% Expanding the preceding boundary expression gives
% E_static-E_0=U_0+2pi G ell^2 int rho_0^2 dV
%                 +(ell^2/2)int rho_0 |grad u_0_vector|^2 dV
%                 +O(G ell^4 with four source derivatives)+O(G^2).
% Integrations use int rho_0 r rho_0' dV=-(3/2)int rho_0^2 dV and
% int rho_0 u_0_vector.Delta u_0_vector dV
%       =2pi G int rho_0^2 dV-int rho_0|grad u_0_vector|^2 dV.
% An independent radial check yields the same correction as
% G ell^2[(9/4)int rho_0 M_0/r^3 dV
%                          +4pi^2 int r^3 rho_0^3/M_0 dV].
% It differs from the once-filter Newton correction pi G ell^2 int rho_0^2.
% Thus the finite-average same-boundary mismatch is a calculated nonzero
% term, not merely an unspecified possible correction. This is a smooth
% slow-static source expansion; microscopic gradient and state errors must
% also be controlled when a fermion state is used to realize the profile.
%
% Reference feedback cannot erase a native response curl on a sufficiently
% rich regular affine covariance branch. Let A be the actual native-source-
% to-native-field response R K W, or its literal-sign counterpart, and Pi
% the symmetric
% spectral susceptibility. With T=(I+A Pi)^(-1)A, exact multiplication gives
% T-T^T=(I+A Pi)^(-1)(A-A^T)(I+Pi A^T)^(-1).
% Indeed multiplication on the left and right cancels the identical
% A Pi A^T terms and leaves A-A^T. No inverse of A itself is required.
% Thus an invertible implicit response cannot turn a nonzero source curl
% into a symmetric response. For a covariance functional gradient the
% antisymmetric bilinear form is b(delta C_1)^T(T-T^T)b(delta C_2).
% If b is surjective on the admitted source tangent space, nonzero A-A^T
% forbids delta E=Tr(h delta C) on that affine chart. An actual branch with
% 0<C<I and linearly independent retained vertices has this surjectivity
% for unrestricted Hermitian covariance tangents (modulo any fixed-charge
% constraint, whose quotient must be taken explicitly). This is a precise
% conditional obstruction for the binding point-field Hartree generator,
% not a theorem against every possible reduced bracket. A fixed-spectrum
% or pure-Slater orbit needs its own detecting tangents or surjectivity;
% a commuting-with-C ambiguity in a dynamical gradient is not discarded.
% The continuum Gaussian test above supplies a nonzero continuum curl.
% Passing it to a finite native parent requires retaining its detecting
% range under the actual source allocation and point readout. A dual
% Galerkin projection containing the two test functions does so exactly;
% a consistent resolved site quadrature does so for sufficiently fine
% sampling by convergence of the two nonzero test pairings. An arbitrary
% finite native map is not assumed faithful merely by suppressing R.
%
% There is also a direct restricted-tangent weak-field witness. Assume a
% regular gapped branch parametrized by lambda=sqrt(G), bounded implicit
% inverse, and h_lambda=h_0+lambda V+O(lambda^2). At lambda=0 take an active
% positive-band subspace Q with occupations separated from 0 and 1, and
% assume Q Delta V Q !=0, where hbar-h=lambda Delta V+O(lambda^2).
% A finite Hermitian tangent delta C=Q Delta V Q is then admissible for
% sufficiently small amplitude; it is not a pure-Slater tangent. Transport
% the active subspace smoothly on the actual branch. Since the field
% variation and hence delta P are O(lambda),
% Tr[(hbar-h)delta C]=lambda ||Q Delta V Q||_HS^2+O(lambda^2),
% Tr[(hbar-h)delta P]=O(lambda^2).
% Thus the exact equality criterion fails on this branch. The nonzero
% compression can be tested with resolved positive-band radial-current
% packets in an exterior interval: Delta v=-9 sigma_ell^2 v/(8r^2)+...
% has a nonzero diagonal current matrix element. The packet and active
% occupation premise must actually be met; no existence theorem for all
% self-consistent branches or all state domains is inferred from this test.
%
% 10. Actual E4 comparison after the independent broader save.
% Read strand_E_operator_mass.tex, its standalone scope and the single
% review_E_operator_mass.tex in full, including every SOURCES entry.
% OM1--OM2 use a different specified positive compact kernel convolved
% with a cell box. The radial integral int_0^1 t^2(1-t^2)^5 dt=256/9009
% verifies their normalization. Our Gaussian kernel is not identified
% with it. Both supply a real cumulative 0<=F<=1 and the same exact
% native vertex A={F,h_0}/2. OM3 agrees with our section 5, including
% the spin trace, factor 1/4 and absence of a further valley multiplier.
% OM4 is an explicit compact-window witness: if L>4(R+b_ell), all nonzero
% weights have cos(2pi x_1/L)>=cos(2pi(R+b_ell)/L)>0. Thus the q=(2pi/L,0,0)
% Fourier coefficient is nonzero. For m>0 and L/a>2 its single p=q,k=0
% term makes Var_0 M strictly positive. This directly validates the domain
% premise used in our positive-square argument without assuming genericity.
% The neutral two-state compression in OM5 is independently checked by
% taking Xi=M Omega/||M Omega||: its matrix is [[0,sigma],[sigma,d]],
% with negative lower eigenvalue. Since ||M||<=Tr|A|, the stated quantitative
% negative upper bound follows by monotonicity of that eigenvalue in d.
% It is a variational compression, not an invariant two-state dynamics.
%
% The source-level mass in OM7 includes the current term and agrees with
% our explicit literal/compatible sign labels (their varsigma=-our sigma).
% Neither the E4 witness nor our witness rules out every ordered solution
% of that coupled charge. OM8 supplies a genuine possible domain check:
% intersect the nonnegative spectral subspaces for a finite list of masses,
% then repeatedly intersect with the preimages under the retained Hermitian
% masses, currents and proposed generator. In finite dimension each strict
% step lowers dimension; the stable space is invariant and, by Hermiticity,
% reducing. Induction proves maximality inside the initial intersection.
% It can be zero and cannot contain the reference vacuum when its mass
% variance is positive. A positive compression without this invariance
% check changes the theory; a successful check still does not solve the
% remaining Gauss row. This is a constructive exact-domain option alongside
% our state-specific positive-Hartree expansion, not an assumed result.
%
% OM9 agrees with our finite-CAR covariance identity for every quasifree C.
% OM10 deliberately multiplies the mean-subtracted mass operators; it is
% NOT the full-reference-Wick quartic convention of sections 3 and 8.
% For V_T=(1/4)sum T_ij {M_i,M_j}, its expectation is
% (1/2)sum T_ij[mu_i mu_j+Re Tr(C A_i A_j-C A_i C A_j)].
% At fixed A,P,T its derivative is
% sum T_ij[mu_j A_i+{A_i,A_j}/4-A_i C A_j], exactly as in E4.
% The one-body contraction may not be omitted while naming this same
% operator. Conversely choosing full reference Wick order changes that
% operator and produces our D=C-P formulas. The ordering distinction is
% explicit; it is not a disagreement in the CAR contractions.
%
% The broader occupation bound can also be checked without a rest-mass
% approximation. For C=P+Gamma, 0<=Gamma<=1-P and N_f=Tr Gamma, decompose
% A into occupied/empty reference blocks. Direct multiplication gives
% Var_C M=Var_P M-Tr[Gamma A_+- A_-+]
%              +Tr[Gamma A_++^2-Gamma A_++ Gamma A_++]
% <=Var_P M+||A_++||^2 N_f.
% If the actual mean obeys mu>=e_* N_f>0, then
% Delta M/mu<=sqrt(Var_P M+||A_++||^2 N_f)/(e_* N_f).
% This agrees with OM11--OM12. Uniform ||A_++||/e_* and
% Var_P M/e_*^2=O(N_f) suffice for the square-root-number upper bound;
% neither follows from N_f alone. General non-Gaussian states also have
% their connected four-point cumulant. Our fixed-number inside/outside
% example already violates a universal decreasing relative fluctuation.
%
% The [X] comparison is restricted to these matched finite-parent covariance,
% variance, number-bound and direct energy-only domain results. Named agreeing
% methods are native projector/interband multiplication and finite-CAR Wick
% contraction, and the positive-square theorem versus neutral-sector
% variational compression. E4 supplies no gravitational coefficient T,
% current-dependent operator reduction or exchange number, so neither the
% weak kernel nor the general K/K_eff receives an external-agreement tag.
% Its adjoint source selector also does not establish equivalence with the
% binding point-field matter generator; the exact criterion in section 9
% must still be checked. The actual E4 review accepts its bounded scope;
% the author's introductory review-pending wording is not promoted into
% an assertion that the outstanding physical reductions have been done.
%
% 11. Actual regular-centre E4 continuation, at its different coupling.
% After our broader derivation and finite-ell static mismatch were saved,
% read strand_E_regular_center_exchange.tex and its scope completely,
% including SOURCES. Its introduction and index say review pending.
% RCX1 uses the compact/cell kernel of OM1 with coordinate variance
% s_ell^2=ell^2/15+a^2/12. RCX2 couples the microscopic static particles
% to V=A^*W, the adjoint-averaged flow, and rebuilds the actual reference
% for h_ell. This is a different field coupling from our binding point
% prescription. The source contract and stationary state realization are
% explicit premises, not supplied fermion solutions.
%
% The conditional boundary coefficient in RCX5 checks directly:
% B_ell=-sum_x mu_x |(A^*W)_x|^2/2. Kinetic energy gives +1/2 and the
% integrated current term gives -1 times this same square. Its variation
% is (A^*nu)_x-|V_x|^2/2: use delta W=delta M_0/(r W),
% nu'=-4pi r J_1/W and integration by parts. Euler homogeneity then gives
% sum mu A^*nu=B_ell; direct matter insertion has coefficient 2 B_ell.
% Thus the half-count distinction agrees with our boundary derivation
% when the state/coupling premises match. Neither source contract alone
% is a solution of the noncommuting constraints.
%
% RCX7's K_xy=int w_r(x)w_r(y)dr/r^2 has TWO full compact/cell filters,
% unlike our one-Gaussian-filter Newton prescription. Its tail is exactly
% 1/R_+ once every weight is one. With c=4pi||k||_infinity/3, integration
% of min(1,c^2 r^6)/r^2 gives K_+=(6/5)c^(1/3); K_-=1/R_+ is valid.
% The quadratic U_ell=-sum K_xy mu_x mu_y/2 and its adjoint Hartree
% derivative are correct. It differs from B_ell by
% (1/2)sum mu_x[(A^*|W|^2)_x-|A^*W_x|^2]>=0.
% Jensen and total kernel variance 3s_ell^2 bound this by
% (3/2)s_ell^2(sum mu_x)Lip(W)^2, as in RCX8. This is an exact positive
% filter-error certificate for THAT adjoint source contract. At zero
% filter width it and our static boundary expression reduce to U_0.
% At finite width it is not the point-coupled expression of section 9.
% The difference can be calculated exactly after choosing the SAME source
% allocation and array for the comparison. Set W_x=R W and V_x=A^*W.
% Our point-static coefficient is
% B_point=(1/2)sum mu_x|W_x|^2-sum mu_x W_x.V_x, whereas
% B_adj=-(1/2)sum mu_x|V_x|^2. Hence
% B_point-B_adj=(1/2)sum mu_x|W_x-V_x|^2>=0.
% For a centered smooth average this difference is O(s_ell^4) in the
% first-G coefficient. Thus the two constructions agree through the
% computed second-moment order, despite different finite-width couplings.
% Indeed the two-filter U has correction 2pi G s_ell^2 int rho_0^2,
% and the Jensen term contributes
% (s_ell^2/2)int rho_0|grad u_0_vector|^2 at leading order.
% These reproduce our independently saved section-9 coefficient after
% matching second moments. This agreement does not identify the two full
% finite-filter kernels or the two source-state contracts.
% RCX9 also keeps the distinct lapse-response error; small energy error
% alone does not bound the Hartree generator error.
%
% RCX10--RCX12 explicitly choose a finite-parent full-reference-Wick
% quantization of U_ell. Its Gaussian Hartree/Fock formulas match our
% Wick trace algebra with their actual two-filter K, not with our k_ell.
% Its non-Gaussian correction is the retained connected four-point
% contraction. The bounded finite Hamiltonian avoids a local mass square
% root by choosing a different, truncated boundary quantization; this
% does not construct the original N,v operator functionals. The stated
% perturbative ground-energy bound follows from the sea trial and the
% complement gap m: if G||V_1||<m/2, completing the square gives
% -2G^2||V_1||^2/m<=E_gs<=0. A scalar ground-energy shift does not make
% the ground vector the original reference or remove other contacts.
% RCX13's rest-particle bounds follow by positivity of the compressed
% E_x and sum_xy Tr(Gamma E_x Gamma E_y)=Tr Gamma^2=N_f, exactly the
% matrix-weight argument of section 4. Their constants retain the actual
% parent/filter dependence. This does not extend them to the full sea.
%
% The actual mathematical agreement covers the conditional unsmoothed
% boundary limit, the explicit filter variance identity, and the stated
% Wick/number-bound algebra. Different finite-filter couplings and the
% unconstructed stationary source family are retained. No external review
% pass, numerical result, universal kernel or complete operator-H2 reduction
% is inferred; the general K/K_eff remains [R] at its declared domain.
%
% 12. Anchored operator annulus: actual source comparison.
% Read the current solve/strand_E_annular_operator.tex, its actual
% review_E_annular_operator.tex and delivered scope, including SOURCES.
% The current author incorporates the review's qualifications; the older
% delivered author did not yet contain all of them. This comparison uses
% the actual reviewed author. AO1--AO6 provide a concrete alternative
% domain construction. On a finite Cartesian parent with all averaging
% collars retained in r_0<r<r_1, first trace the c-number reference
% r_A(f)=Tr P(h(f))H_A, then Weyl-order its Taylor series in Hermitian
% Fock selectors. This added quantization is not an operator-valued
% spectral projector traced over one-body space. The resolvent expansion
% converges if 2(sum_A||H_A||)max_A||F_A||/m<1. A fixed external charge
% M_b>sup_r||D_0(r)|| makes M_0=M_b I+D_0 positive on the FULL finite
% Fock space. Integrating the ordered Gauss rows gives at t=sqrt(G)=0
% the derivative [[0,R],[0,0]], so I-DPhi is invertible. The analytic
% implicit-function theorem gives a local branch with M>=delta I/2 at
% fixed parent, anchor and annulus. This checks independently by the
% triangular linearization and positive-root functional calculus. It
% neither removes the anchor nor supplies a uniform isolated-object limit.
% Spatial C^5 regularity and analyticity in t are correctly distinguished.
%
% AO7--AO9 retain the effects independently identified in sections 7--8:
% B_1=-varsigma int dV Sym(W_0,J_0),
% B_2=-T(F_1)-varsigma int dV[Sym(W_1,J_0)+Sym(W_0,J_1)],
% T=sum_occ,empty (Z^dag Z+ZZ^dag)/[2(E_b-E_o)]>=0.
% Tr h_0 DP_0[V]=0 removes the integrated first reference variation,
% not the quadratic G contact. Source variation of B_1 has derivatives
% L_1 and -varsigma W_0, matching our compatible-sign result only at
% varsigma=-1. Its selectors are adjoint averages, a different coupling
% from the binding point fields. Direct insertion also has the lapse-work
% term int dV Sym(L_1,rho_0); the reviewed author leaves its cancellation
% and optional inner-clock work unresolved. No insertion identity follows.
%
% AO11--AO13 expand about dominant M_b I to obtain an anchored
% mass--current pair of order sqrt(G)/sqrt(M_b). For its MEAN-SUBTRACTED
% PRODUCT ordering, Gaussian contraction gives
% mu nu+Re Tr C e(1-C)p, and differentiation gives
% nu e+mu p+{e,p}/2-eCp-pCe. These agree with our finite-matrix
% contraction/variation methods and section 10's distinct product ordering.
% They are not our full-reference-Wick density pair or an isolated first-G
% kernel. Its bounded outer charge can separately be chosen as a unitary
% number-preserving generator; insertion equality and H2 closure remain
% unproved. This is a constructive domain option with additional boundary
% input, not an assumed resolution of the direct energy-only obstruction.
%
% 13. Actual numerical task-8 controls: producer and observable audit.
% Read SPEC_final4b_engine.md section 11.6 and complete local producers
% engine_hf_cluster_v1.py, engine_hf_cluster_v2.py and
% engine_massfluct_v1.py in .wt-solve_a/solve/code/engine.
% Read the local hf_cluster_smoke_ns7.json and actual remote
% <internal>/hf_cluster_v2_scan1.md, its JSON
% metadata/representative rows/convergence flags and completion log.
% Executed remote v2 and its imported v1 are byte-identical to the local
% sources read. The smoke names v1; the scan names v2 importing v1.
% This audit did not execute a scientific job or reconstruct an output.
%
% These controls use ONE two-component kappa/r radial chain, vertices
% A_n={Pi_n,h_chain}/2 and V_nm=-G/max(R_n,R_m), including n=m.
% They have no source smoothing, exact native Cartesian decomposition,
% or actual current/reference-response kernel. The scan covers n_s=6,7,8,
% kappa=1,2, N_f=1,...,4 and targets
% |E_pot|/(Tr h C_m-N_f m)=.01,.03,.1,.3 at N_f=4. Its ratios are
% finite controls of this chosen chain Hamiltonian, not measurements of
% the declared four-component cubic theory or a large-N exponent.
%
% Each N_f sector solves a Hartree sea P_N, then freezes it in
% H_sym(P_N)=dGamma(h)+1/2 sum_nm V_nm
%   [dGamma(A_n)-Tr P_N A_n][dGamma(A_m)-Tr P_N A_m].
% Its vacuum subtraction uses P_0: different referenced Hamiltonians,
% not two sectors of one fixed operator. The first-order label E_x is
% the FULL connected product contraction
% 1/2 sum V Tr C A_n(1-C)A_m, with a positive-band E_x_matter split.
% It includes a one-body contraction, unlike section 3's full-reference-
% Wick Fock term. The covariance matrix is positive semidefinite, but
% its entries need not be nonnegative; the negative-semidefinite monopole
% kernel makes the total contraction nonpositive without that assumption.
%
% The executed 'wick' option omits a reference-dependent scalar. Put
% Sigma_P=1/2 sum V A_n(1-2P)A_m and
% F_P=1/2 sum V Tr[P A_n(1-P)A_m]. The reference-Wick interaction is
% H_sym(P)-dGamma(Sigma_P)+Tr(P Sigma_P)-F_P. The code subtracts only
% dGamma(Sigma_P). Its missing scalar is
% Tr(P Sigma_P)-F_P=-1/2 sum V Tr(P A_n P A_m),
% which need not cancel for different P_N. The reported ordering shift
% is the shift of the executed operators, not the normalized full-Wick
% comparison. v2 correctly suppresses symmetric-HF comparisons for this
% option, but that does not restore the scalar. A corrected comparison
% must retain it and specify a common reference/operator before comparing
% sectors. No such corrected run is inferred. The existing RED derived
% this same scalar independently from the Gaussian expansion in C-P.
%
% Self-consistent HF updates its reference sea, while the exact operator
% freezes the Hartree sea. Their proximity alone is not a controlled
% higher-order comparison of one Hamiltonian. Actual JSON records several
% hf_sc_converged=false cases at 2000 iterations and target .3; these are
% not converged HF results. Small-coupling values remain raw chain controls
% under their stated definitions. Four-point fits at N_f=1,...,4 prove
% no universal relative error: for example n_s=8,kappa=1,target .01 gives
% fitted power 1.049 for E_x_matter/E_pot and -0.891 for the collar matter
% fluctuation. Neither is the positive-particle bound or a universal law.
%
% The printed fluct_exact is sqrt(max(Var_N-Var_vac,0)) divided by a
% vacuum-subtracted mean: clipped EXCESS variance. Mean-source subtraction
% leaves Var_N itself unchanged. Its strong-scan zeros are not vanishing
% quantum variance; the v1 smoke retains full variance profiles separately.
% engine_massfluct_v1 uses sharp cumulative chain energies without ell
% weights or the coupled linear-response A_eff. Its factor 16 kappa is
% conditional addition of asserted independent chain copies, not an exact
% native cubic decomposition supplied here. The channel stop follows
% positive matter occupation and does not itself bound the remaining sea
% variance. No full-model [N] or numerical [X] follows. A matched check
% still requires identical native vertices, filters, ordering, reference,
% current-dependent mass response and state/domain on both sides.
%
% 13.1. Actual v3 Wick scan: corrected pair contraction, unequal references.
% Read on compute-host, without running or copying its calculation:
% <internal>/engine_hf_cluster_v3.py,
% its imported v1/v2 definitions, and engine/HF/
% hf_cluster_v3_wick_scan1.{json,md,log},
% hf_cluster_v3_wick_scan2.{json,log}. The records name v3 and its logs
% contain the matching rows. This identifies the available producer and
% results, not an unrecorded historical byte identity or unfinished job.
% The v3 ex_wick and sigma_wick now implement exactly
% E_F(D)=-1/2 sum_ij V_ij Tr(A_i D A_j D),
% Sigma_F(D)=-sum_ij V_ij A_i D A_j.
% At rank-one D their pair exchange cancels the Hartree self-energy;
% every N_f=1 Wick row accordingly has y=E_F/E_H,pair=-1. This corrects
% the earlier mean-subtracted-product contraction for this calculation.
%
% Its exact_with_state still evaluates
% H_raw(P)=dGamma(h)+H_sym(P)-dGamma(Sigma_P),
% Sigma_P=1/2 sum_ij V_ij A_i(1-2P)A_j,
% H_sym(P)=1/2 sum_ij V_ij[dGamma(A_i)-Tr(P A_i)]
%                              [dGamma(A_j)-Tr(P A_j)].
% The symmetric sum makes the product Hermitian. The desired full
% reference-Wick pair differs by the scalar
% c(P)=Tr(P Sigma_P)-F(P)
%     =-1/2 sum_ij V_ij Tr(P A_i P A_j),
% F(P)=1/2 sum_ij V_ij Tr[A_i P A_j(1-P)].
% In addition, the declared centered kinetic term is
% dGamma(h)-Tr(hP). Therefore the complete fixed-P correction is
% k(P)=-Tr(hP)+c(P), and a sector difference needs k(P_N)-k(P_0).
% run_case actually passes the distinct sector Hartree projectors H[C0]
% and H0[C0] to exact_with_state. The producer header's assertion that
% the omitted scalar is sector-independent is not valid for those inputs.
% Neither constant generally cancels between these distinct references.
% This is an exact finite-CAR algebra finding, not a numerical error fit.
%
% A scalar correction leaves each fixed-sector eigenvector, its covariance
% and its mode-cut entropy unchanged. Thus the available exact SVD of the
% inner 2 floor(n_s/2) modes is a legitimate bipartition observable of the
% specified finite chain Hamiltonian, unaffected by this energy-zero repair.
% It is not a three-dimensional geometric cut. Nor does a scalar repair
% make independently updated Hartree/HF seas and the exact frozen Hartree
% sea one common Hamiltonian or the binding transported reference.
% hf_wick_sc computes convergence and iteration data internally, but v3
% does not retain those fields in its output rows; convergence is not
% inferred from a printed energy. The n_s=6,7,8, kappa=1,2 scan at targets
% .01,.03,.1 is a finite two-component-chain control. The available scan2
% has n_s=10,kappa=1 at targets .03,.1 with N_f=1,...,8. Its free fitted
% powers are .825205 and .805683; fixed-power-one maximum relative
% residuals are .177697 and .192987, versus .595427 and .571836 at power
% one-half. These finitely many points prove no asymptotic exponent.
% At n_s=10,target .1,N_f=8 the raw (exact-HFsc)/E_pot is .00697989;
% it supplies no universal 1e-5 accuracy, even before the reference repair.
% The requested n_s=12 is not a result in the read scan2 record.
%
% The executed chain still has two components, angular kappa/r, and
% V(r,s)=-G/max(r,s), with no declared ell source map or current-dependent
% Gauss charge. Its a=1 parameter G is g, its mass is ma=.5; neither
% sector statistics nor a mode-cut entropy restores the missing native
% cubic operator/cut. Thus no full-model numerical acceptance follows.
% Actual local SPEC_final4b_engine.md section 12.7 and observables section
% 8.8 were read against these producers. Ordering is a named prescription,
% not a result of the binding constraint reduction. The actual engine
% section 13.1 reference substitution is checked separately below and in
% declared_local_area.tex; it is not silently used in this note.
%
% 14. Declared transported sea: constructive Hartree state space [R].
% The current full binding GOAL specifies the incoming sea transported by
% the model's evolution, never instantaneous band refilling. At the common
% one-body Hartree level, let P_T=U P_in U^dag, C=U C_in U^dag and D=C-P_T.
% Then exactly
% i dot D=[h(D),D],  q_I=Tr(B_I D).
% A positive added-particle initial D=Gamma remains positive with its
% spectrum and N_f=Tr Gamma fixed. This positivity of the covariance
% difference does not make a local energy vertex or enclosed energy
% positive: that is still a separate state/domain question.
%
% On the smooth compatible-sign unsmoothed branch of section 7,
% delta E=Tr(h delta D). Thus E[D], on the coadjoint orbit of D, gives
% exactly the displayed source evolution with the Lie--Poisson bracket
% -i Tr D[grad F,grad G] on that Hermitian-matrix orbit. Along it
% dot E=-i Tr h[h,D]=0. This autonomous D-source description requires no
% instantaneous Pi pullback. It is not yet the native CAR bracket of
% affine reference-subtracted observables: at fixed reference R=P_T,
% that bracket is -i Tr C[A,B]=-i Tr D[A,B]-i Tr R[A,B].
% A constructive native Hartree history instead uses, for every fixed R,
% E_R(C)=E_boundary[C-R], whose fixed-R gradient is h. The ordinary
% C-covariance bracket gives i dot C=[h,C], with the separately declared
% companion i dot R=[h,R]. Their difference is the D equation and
% dE/dt=-i Tr h[h,C-R]=0. This does not construct an autonomous C-only
% Hamiltonian after the transported history is eliminated. A product
% orbit with an opposite mathematical bracket for R also realizes the
% D bracket, but no extra physical reference degrees of freedom or
% independent Hilbert-space factor are declared by that representation.
% Its source Hessian is K of
% section 7 in affine D coordinates, while an orbit's curved paths also
% retain Tr h delta^2 D. Its native sqrt(G) and O(G) current kernels are
% unchanged at fixed D. It is not a constructed state-independent
% full-Fock quantization of nonlinear operator constraints.
%
% At finite averaging the boundary gradient in D is hbar, containing
% adjoint selectors, whereas the binding evolution uses point-field h.
% Equality on all allowed D tangents requires Tr[(hbar-h)delta D]=0;
% on the D orbit a commuting-D difference is invisible to source evolution.
% Equality of the physical C and reference R evolutions separately needs
% both [hbar-h,C]=0 and [hbar-h,R]=0, a stronger condition. A universal
% equality or a response-curl obstruction still needs the actual native
% source/readout maps and detecting tangents. The source curl in section 9
% persists without any reference feedback; its static spectral Pi
% congruence identity is not needed for this D-state calculation. Replacing
% h by hbar is an alternative coupling, not a silent binding-model repair.
%
% If physical C rather than D is held as the variation coordinate while
% the history changes, the reference response is causal. For fixed P_in,
% delta P_T(t)=-i int_{t_i}^t U(t,s)[delta h(s),P_T(s)]U(t,s)^dag ds.
% Add U delta P_in U^dag when the initial reference also varies. Therefore
% Pi^R_IJ(t,s)=-i Theta(t-s)Tr[B_I U(t,s)[B_J,P_T(s)]U(t,s)^dag]
% is a retarded susceptibility, not the real symmetric spectral Pi.
% With actual A=R K W, delta q=b(delta C)-Pi^R*delta phi and
% delta phi=A delta q form a Volterra response problem, subject to its
% own existence and inverse bounds. The exact arbitrary endpoint variation
% is delta E=Tr[h(delta C-delta P_T)] without averaging, or
% Tr[hbar(delta C-delta P_T)] with averaging. In general Tr h delta P_T
% does not vanish. Along the common time evolution Tr h dot P_T=0 does
% vanish, a different statement. No spectral contact cancellation or
% symmetric static inverse is silently transferred to this history.
%
% The local-area note gives an explicit positive-lapse, zero-flow smooth
% pulse returning to h_0 but leaving P_T different from P_0. Thus the
% no-pocket threshold alone does not identify these references. Uniform
% spatial fields are a special exact commuting case. A gapped temporal
% adiabatic limit can connect them only with quantitative state/history
% control. The instantaneous-reference E4/RCX/AO comparisons above retain
% their own scope and do not establish the transported operator theory.
%
% For any prescribed transported Slater reference, the centered bilinear
% mass dGamma(A)-Tr(P_T A) still has variance Tr P_T A(1-P_T)A. If that
% is positive, its mean-zero spectral obstruction remains exact. The
% explicit Fourier formula in section 5 is its flat P_0 specialization.
% Gaussian full-reference Wick algebra likewise works at each fixed P_T
% if an ordering is explicitly frozen there; the first-G P_0 prescription
% has the fixed-history perturbative scope stated in section 3. Beyond
% a common quadratic Hartree history the transported reference can be
% non-Gaussian. Its higher connected contractions must then be included;
% it is not determined by a one-body projector and Wick's theorem alone.
%
% Actual further comparison [X, transported algebra only]: fully read
% .wt-solve_b/solve/strand_B_causal_sea_budget.tex, its review and
% staged scope. CS1--5 separately derives common-unitary C,P_T,D transport,
% exact propagated reference mismatch and trace-duality source bounds.
% CS12--14 derives regional reference-subtracted matter work and flux.
% These agree with the direct Duhamel/covariance calculations here and
% the explicit detector/source distinction in declared_entropy.tex.
% That source uses an adjoint-field prescribed-background preparation;
% it supplies neither the binding point-field reduction nor an operator
% current-dependent interaction. Only the named common-unitary algebra
% and error/flux identities are an actual cross-lane agreement.
%
% 15. Regular-centre quantum flux: an exact operator-domain obstruction [R].
% Independently derived and saved before the existing RED's new duty-9
% continuation was read. This strengthens the leading energy-only mass
% test in section 5 without replacing the current-dependent constraint by
% that approximation. It concerns a precise regular-centre ordering class,
% not every possible quantization with extra quantum flux terms.
%
% Specify the domain first. On a finite native fermion parent the Fock
% space is finite dimensional. Alternatively use the incoming-sea CAR
% representation and a common invariant domain containing its vacuum.
% The theorem assumes a regular spherical Hermitian extrinsic curvature,
% K_ij(0)=delta_ij k with self-adjoint k. It does not require the clock
% or the matter source to be a positive operator. For example, regular
% fields N(x)=N_c+O(r^2), v_i(x)=x_i B+O(r^3), with self-adjoint
% invertible N_c and self-adjoint B, admit the declared Hermitian order
% K_ij={N^(-1),S_ij}/2, S_ij=(partial_i v_j+partial_j v_i)/2.
% Then k={N_c^(-1),B}/2. If a positive invertible N is available on a
% chosen domain, K_ij=N^(-1/2) S_ij N^(-1/2) is another explicit
% Hermitian ordering with the same centre property. Declare the literal quadratic
% Einstein flux in these K operators, with no additional commutator or
% subtraction term in the geometric flux. Then exactly
% rhohat_ell(0)=(1/(16pi G))[(sum_i K_ii)^2-sum_ij K_ij K_ij]
%             =3 k^2/(8pi G) >= 0.                         (QC1)
% Any permutation or Hermitian symmetrization of the two identical
% centre k factors gives the same square. The sign in the momentum row
% and its nonzero current terms do not change this centre identity.
% No expansion in G or small current has been used.
%
% The declared compactly averaged fermion source is instead
% O=dGamma(A_w)-Tr(P_T A_w), A_w={W,h_0}/2,
% where W_xx=w_x I_4, w_x>=0 are nonzero finite-support density weights.
% At the incoming flat sea P_T=P_0 its mean is zero and its exact
% physical variance is sigma^2=Tr[P_0 A_w(1-P_0)A_w].
% The flat-sea sign test also applies to a nonzero summable nonnegative
% window, including a Gaussian: A_w is trace class, its Fourier weight
% is continuous and nonzero near zero, and the interband variance integral
% is strictly positive. For m>0 take nearby p!=0,q=0; for m=0 take
% p=(epsilon,0,0),q=(-2epsilon,0,0), epsilon>0 sufficiently small.
% Unequal energies and distinct normalized Dirac symbols persist on
% open neighborhoods, giving a positive-measure integrand. This Gaussian
% extension is an incoming-flat-sea statement; its value at a later
% history must be checked, not inferred from compact-support faithfulness.
% The strict finite-region faithfulness proved in
% declared_faithful_vacuum.tex makes sigma^2>0 for this nonzero local
% centered bilinear. It also proves that O has both positive and
% negative spectral values. The same statement holds after every finite
% bounded finite-range quadratic history, by the transported-faithfulness
% theorem in that source. This is not an instantaneous-refilling claim.
% Thus (QC1) cannot hold on the full native Hilbert space including this
% reference. Equivalently, its zero reference expectation would force
% k|Omega_T>=0 and hence O|Omega_T>=0, contrary to sigma^2>0.
% For an unbounded k these statements require the common domain to
% contain |Omega_T> and the quadratic-form identity; the finite-parent
% statement below has no unbounded-domain issue.
%
% There is a concrete finite-parent witness independent of the infinite-
% parent faithfulness proof. On a periodic n^3 parent of side L=na,
% let Wtilde(k)=n^(-3)sum_x w_x exp(-ik.a x), E_p=sqrt(m^2+sum d_i(p)^2),
% d_i(p)=sin(ap_i)/a, and fill the negative-energy bands. Direct spin
% tracing of P_p A_pq(1-P_q)A_qp gives
% sigma^2=(1/4)sum_pq |Wtilde(p-q)|^2 (E_p-E_q)^2
%           [1-(m^2+d(p).d(q))/(E_p E_q)].                 (QC2)
% This equality is for m>0. For m=0, its sum over E_p E_q>0 is a
% lower bound; the remaining zero-mode terms, evaluated with the specified
% filling rather than this divided expression, are nonnegative. The
% witnesses use only nonzero-energy momenta. Suppose all w_x lie in |a x_1|<=b
% in a coordinate chart not crossing the periodic seam.
% For m>0 choose q=0,p=(2pi/L,0,0). If L>4b and n>2,
% Re Wtilde(p)>0, E_p>m, and the bracket is 1-m/E_p>0.
% For m=0 take delta=2pi/n, n>6,
% p=(delta/a,0,0), q=(-2delta/a,0,0). The two energies
% sin(delta)/a and sin(2delta)/a are positive and unequal, while the
% bracket is 2. If L>12b then Re Wtilde(p-q)>0. Consequently the
% reference has positive variance even on this finite native parent.
% The compact weights can be the actual smooth ell-average sampled on
% the lattice; this witness requires a parent large enough to contain
% that window and does not use a shell-degeneracy replacement.
%
% The obstruction has an operator-norm version. Set
% |Xi>=O|Omega>/sigma, d=<Xi|O|Xi>. These two unit vectors are
% orthogonal and have the same fermion number. The compression of O is
% [[0,sigma],[sigma,d]], whose lower eigenvalue is
% lambda_-=(d-sqrt(d^2+4sigma^2))/2<0. For every positive operator B_+
% on a space containing these two vectors,
% ||O-B_+|| >= (sqrt(d^2+4sigma^2)-d)/2
%            >= (sqrt(B_O^2+4sigma^2)-B_O)/2               (QC3)
% whenever B_O>=||O||. Therefore a positive geometric square cannot
% approximate this full source in operator norm arbitrarily well at
% fixed parent, a, ell and reference. This is a domain statement, not a
% numerical lower bound independent of the averaging volume.
%
% Ordering qualification and the regular weak-g limit.
% The square argument must not be extended to every ordering of N and
% B separately. An explicit counterexample is X=diag(1,e), 0<e<1,
% Y=[[1,1],[1,1]], N=X^(-1/2), B=Y^(1/2). Both N and B are
% Hermitian and N is positive, but the equally classical centre monomial
% ordered as (N^(-2)B^2+B^2 N^(-2))/2 has matrix
% [[1,(1+e)/2],[(1+e)/2,e]], of determinant -(1-e)^2/4.
% It is indefinite. This ordering changes the quantum flux by commutator
% terms relative to the Hermitian-curvature square. The example evades
% that square test only; it supplies neither the momentum constraint,
% Hartree reproduction, a full radial domain nor the required reference.
%
% There is nevertheless a precise perturbative obstruction for a wider
% ordering class. At fixed finite parent and fixed ma,ell/a, suppose
% A_g=N_c(g)^(-1) obeys epsilon_g=||A_g-I||->0, and
% X_g=a B(g)/sqrt(g) is Hermitian with ||X_g||<=C uniformly.
% Take any fixed finite real linear combination of Hermitian-symmetrized
% words containing two A_g and two X_g, with coefficients c_j summing
% to 1 so that its commuting limit is A_g^2 X_g^2. Let C_ord=sum|c_j|.
% Replacing the two A_g factors in each word by I gives the norm bound
% ||rhotilde_flux-(3/(8pi))X_g^2||
%   <=(3/(8pi)) C_ord epsilon_g(2+epsilon_g) C^2.           (QC3a)
% This includes the indefinite ordering example, but excludes a clock
% remaining an order-one distance from the free clock as g->0, singular
% coefficients, or extra quantum terms not obtained by such factor order.
% Suppose also that the actual transported reference tends to P_0.
% The explicit condition int_0^T ||h_g(t)-h_0||dt->0 at fixed finite T
% is sufficient by Duhamel's formula. Then rhotilde_g tends in
% norm to the nonpositive source rhotilde_0 because reference subtraction
% changes it only by Tr[(P_T(g)-P_0)a^4 A_w] times the identity.
% Let delta_0>0 be the negative-witness bound (QC3), in lattice units.
% Exact operator equality would imply
% delta_0 <= ||rhotilde_g-rhotilde_0||
%          +(3/(8pi)) C_ord epsilon_g(2+epsilon_g)C^2,
% impossible for sufficiently small g. Thus a regular perturbative
% full-space weak-field operator lift in this ordering class does not
% exist. No uniform conclusion is asserted when the parent, state,
% observation time or clock norm scales singularly with g.
%
% For the particular raw order K=N^(-1)B, its Hermitian-square flux is
% Re(K^2)=H^2-L^2, H=(K+K^dag)/2, L=(K-K^dag)/(2i).
% In the same scaled variables ||a L/sqrt(g)||<=epsilon_g C,
% so its negative part is bounded by (3/(8pi))epsilon_g^2 C^2.
% This makes explicit both the finite-g loophole and its regular
% weak-g limitation. It is not permissible to call Re(K^2) positive.
%
% Allowed restricted construction and its actual limits.
% In a specified state with mean rho_bar>0 and variance V, define the
% spectral projection Q_epsilon=1_{|O-rho_bar|<=epsilon rho_bar}, 0<epsilon<1.
% On its range the self-adjoint centre root
% k_Q=zeta sqrt(8pi G/3) sqrt(Q_epsilon O Q_epsilon)
% exists, where zeta=+1 or -1 is a fixed inward/outward branch choice.
% The convergent binomial expansion about rho_bar has small parameter
% ||Q_epsilon(O-rho_bar)Q_epsilon||/rho_bar<=epsilon.
% This constructs a single-centre compressed root, not the full radial
% constraint. The physical spectral probabilities obey
% Prob(O<0)<=V/(V+rho_bar^2) (the one-sided variance bound) and
% Prob(|O-rho_bar|>epsilon rho_bar)<=V/(epsilon^2 rho_bar^2).
% A variance bound controls these probabilities, not exact support.
% In Gaussian states V=Tr[C A_w(1-C)A_w] retains the reference-sea
% contribution. The reached-state estimates in the faithful-vacuum
% note can bound this V and mean for a specified preparation and time;
% they do not supply a universal N_f^(-1/2) law.
%
% Different radii use noncommuting bilinear masses. A collection of
% positive compressions therefore needs an actual common domain,
% compatibility with both constraints, and invariance under the reduced
% Hamiltonian. Compressing a Hamiltonian to that domain generally changes
% its evolution unless the domain reduces it. Excluding the reference
% vacuum avoids the zero-mean contradiction only by changing the claimed
% full-space domain; a positive Hartree expectation does not do this.
% An anchored annulus avoids the regular-centre hypothesis but has the
% different boundary/domain already compared in section 12. Allowing
% nonspherical quantum shear changes the centre flux from a square, but
% is outside the stated spherical geometry. Adding a geometric
% subtraction/commutator term can evade (QC1); its coefficient, ordering,
% constraint consistency and Hartree limit would then have to be derived.
% Merely permuting the two centre k factors cannot supply such a term.
% There is also an exact invariance test, independently supplied by the
% original RED and checked here. In a finite massive parent's fixed-charge
% sector, the free Hamiltonian has the unique filled-sea ground vector
% |Omega>. For a centered local source O with positive vacuum variance,
% its vacuum spectral probabilities on both O>0 and O<0 are nonzero.
% Let Q_+=1_{O>0}. If [H_0,Q_+]=0, uniqueness forces Q_+|Omega>
% to be either 0 or |Omega>, contradicting those two probabilities.
% Thus this positive-source spectral subspace is not invariant even
% under the free dynamics. The same argument applies to 1_{O>=0}.
% This does not exclude every invariant positive subspace that could be
% separately constructed; it rules out adopting the full natural
% positive spectral projection as an unchanged dynamical sector.
%
% First-G pair status. A finite-parent, fully reference-Wick-ordered
% symmetric pair Hamiltonian is a legitimate self-adjoint operator,
% as constructed in sections 3 and 8. Its background Hartree derivative
% and Fock kernel can be checked against an actual reciprocal response
% map, and its physical covariance can be computed. It does not satisfy
% (QC1) on the full native Fock space simply because its Hartree field
% solves the positive-expectation constraints. The native sqrt(G)
% density/current expansion, finite-ell point/adjoint mismatch, and
% transported-history dependence remain separate requirements. No
% first-G universal pair kernel or full-Fock mass-functional geometry
% follows from the positive local-history example alone.
%
% A concrete transported-background quadratic construction makes the
% available approximation explicit. On an actual compatible-sign branch
% where the native source/field map is reciprocal, fix a time, P_T and
% D_bar. Let q_I=Tr(B_I D), qbar_I=Tr(B_I D_bar), phi=grad E(qbar),
% and let K be its symmetric Hessian pulled through the actual source
% and field maps. In the unsmoothed source variables K is the closed
% density/current kernel of section 7. This D-coordinate construction
% uses no instantaneous spectral Pi. At a finite parent declare
% H_quad=c+sum_I[phi_I-sum_J K_IJ qbar_J] qhat_I
%          +(1/2)sum_IJ K_IJ :qhat_I qhat_J:_(P_T),         (QC4)
% with qhat_I=dGamma(B_I)-Tr(P_T B_I), full reference Wick ordering
% and Hermitian symmetrization. Choose c to match the background Hartree
% energy. This is a bounded self-adjoint native-Fock operator. Its Hartree
% derivative is exactly sum_I[phi_I+sum_J K_IJ(q_J-qbar_J)]B_I,
% the linearized declared field on that reciprocal branch. The Gaussian
% Fock correction and its one-body derivative are
% E_F=-(1/2)sum_IJ K_IJ Tr(B_I D B_J D),
% h_F=-sum_IJ K_IJ B_I D B_J.                              (QC5)
% No change of reference mean removes the physical variance. The
% linearized current-dependent mass has vertex
% A_r^D=sum_I (L_r W)_I B_I,
% where L_r is section 7's mass response and W is the actual averaging
% map (identity in unsmoothed coordinates). Its exact Gaussian variance
% is Tr[C_bar A_r^D(1-C_bar)A_r^D]. These expressions use the native
% energy/current matrices and the selected transported background.
% They are a local fluctuation construction, not a quantization of the
% nonlinear constraints on their full spectrum. A nonsymmetric point-
% field response cannot be used as K in (QC4) while retaining its claimed
% Hartree derivative; section 9 states the relevant admissibility tests.
% Choosing a new background or changing its transport changes (QC4).
%
% There is a further algebraic obstruction to a universal two-body
% replacement in the fixed normal-source variables. On any admissible
% scaling family D_t=t D_*, 0<t<t_max, with common incoming reference,
% nonzero I_*=int dV z_0[rho_*] J_* and positive regular M_0, section 7
% gives E_(sqrt G)[D_t]=t^(3/2) I_* (up to the momentum-sign choice).
% The Hartree expectation of a fixed one-body plus fully reference-Wick-
% ordered two-body operator is instead t L+t^2 Q+c. It cannot equal
% t^(3/2) I_* on an open t interval. The statement is conditional on
% such a family actually lying in the allowed state/source domain; it
% does not identify arbitrary source functions with native covariances.
% Even the O(G) current kernel depends on the shape of rho through
% sqrt[M_0(min)/M_0(max)]. A universal background-independent current
% pair therefore cannot reproduce its Hessian over independent density
% profiles. This excludes an exact two-body replacement in these
% coordinates, not a controlled large-occupation approximation or the
% weak static canonical-variable conversion in section 1. Replacing
% expectation of an operator square root by square root of its mean
% also needs the fluctuation/domain estimates above; it is not an
% identity furnished by an ordering prescription.
%
% Lattice units make the domain obstruction explicit at every g=G/a^2>0:
% rhotilde=a^4 O=3 ktilde^2/(8pi g), ktilde=a k.
% The dimensionless variance is a^8 sigma^2 and the dimensionless
% Fourier expression depends on ma, ell/a, the actual window and parent.
% In finite cell-charge coordinates qtilde_I=a q_I, the quadratic
% construction uses Ktilde_IJ=K_IJ/a and a H_quad; its response has
% the stated sqrt(g) and g terms at fixed dimensionless normal sources.
% The spectral-probability bounds depend on V/rho_bar^2. Their smallness,
% and any independence of ell in a<<ell<<geometric lengths, must be
% checked in the actual state rather than inferred from g or N_f alone.
%
% SOURCES.
% GOAL_alt.md full current text read after the hydrodynamic/averaging
% and duties 8--9 update. Required THINK sections and pinned REQUIREMENTS
% already read; duty 9's operator, exchange and fluctuation assertions are
% being checked rather than used as conclusions. No E4 counterpart was read
% before the independent weak-sector and broader-continuation saves.
% No old interaction or geometric construction
% is imported. The matched numerical task 8 comparison remains open;
% actual raw chain controls are now assessed in section 13.
% Actual redteam_declared_exchange_kernel.tex and its SOURCES read in full
% after both independent derivations were saved. Agreed named methods:
% spherical shell/Poisson integration and Hartree functional variation;
% finite-matrix Wick contractions and positive-projector bounds; exact
% reference covariance and the positive-square spectral test. The one-
% versus two-filter distinction was checked by Gaussian convolution and
% the independent radial mass-square factorization. Mandatory tail/domain
% repairs were applied. This same-lane agreement is [R], not manufactured
% external [X]. The subsequent actual E4 source/scope/review comparison is
% recorded in section 10 at its precise [X] covariance/domain scope.
% Actual declared_h2_continuum.tex and its SOURCES were read in full after
% the constructive continuation and again after the admissible-tangent
% clarification. Its author independently checked section 7's Hessian by
% integration by parts. The present author derived the four response blocks,
% reference susceptibility and quadratic ordering independently; the existing
% RED has checked this continuation, including the literal/coherent sign,
% the native response-map qualification, and finite-ell static distinction.
% Its full current source and SOURCES were read after its independent
% derivation and actual E4 comparison. No additional reviewer or gate is used.
% The subsequent actual regular-centre E4 source/scope comparison is in
% section 11; its different adjoint coupling and pending external review
% remain explicit. Its title is not used as evidence of a complete kernel.
% Actual annular source/review/scope and numerical producers/outputs were
% read as in sections 12--13. Their distinct boundary, coupling, ordering
% and observable definitions remain explicit. The literal-sign,
% finite-average operator reduction and full matched numerical comparison
% remain [S].
% Current transported-reference GOAL read fully. Section 14 was derived
% directly from common-unitary covariance evolution, the already derived
% boundary variation, and the Duhamel response. The existing RED separately
% obtained the D-orbit Hamiltonian route and retarded susceptibility. The
% spectral comparisons in sections 8--12 are not binding reference laws.
% Section 15 uses the declared Einstein flux itself at a regular centre,
% exact finite-CAR spectral/variance algebra, and the already independently
% proved faithful-vacuum theorem. It was saved before intake of the same
% RED's new full-quantum continuation. The fresh regular-centre ordering
% obstruction is [R]; no external matching theorem has been read or
% asserted. Its square ordering/domain and allowed alternatives are part
% of the statement, not hidden operator-positivity assumptions.
% The current original RED's complete source and SOURCES were read after
% its independent section 16 save. The centre contraction, neutral-sector
% spectral-distance test, regular weak-g obstruction, tangent full-Wick
% pair and physical linearized variance agree by the explicitly named
% methods: direct Einstein-flux contraction versus the variance/spectral
% theorem; finite-word norm comparison versus the enclosed-mass weak-limit
% argument; and finite-CAR Wick contraction versus functional variation.
% The RED's free-H_0 positive-spectral-sector noninvariance proof was
% independently checked and incorporated with its unique-vacuum premise.
% Its requested Duhamel convergence condition is now explicit, rather
% than inferred from boundedness of a history. These new operator results
% remain [R]; the prior actual OM [X] is only its stated variance/domain
% scope and is not promoted to an external proof of the centre extension.

% 16. UNPROJECTED SPHERICAL OPERATOR CONSTRAINTS AND ANGULAR VARIANCE [R].
% Independent substantive continuation, saved before a new RED or
% mechanics-side continuation was read. The compact statement was unchanged
% at that independent save. This tests the full quantum requirement separately from
% Hartree source tangency: an expectation-level Gauss law is NOT assumed
% to be an operator identity. No stationary state, angular projection,
% alternate reference, new field, code or numerical run is introduced.

% QV1. Exact native incoming-sea variance for every localizable weight.
% Work on the infinite cubic parent, a=1, finite mu=ma>=0 and g=G/a^2>0.
% Write d(k)=(mu,sin k_1,sin k_2,sin k_3), E(k)=|d(k)| and
% P0(k)=(I-h0(k)/E(k))/2, with isolated zero modes irrelevant to the
% infinite-parent covariance. Let w be a nonzero real l1 site sequence,
% W its scalar multiplication operator, and A_w={W,h0}/2. Since h0 is
% bounded and W is trace class, A_w is trace class. The fermionic
% normal-ordered observable
%        O_w=dGamma(A_w)-Tr(P0 A_w)
% is a bounded self-adjoint CAR observable, defined also by norm limits
% of finite-support weights. It has zero incoming-sea mean. Its exact
% physical variance is
% V_mu[w]=Tr[P0 A_w(I-P0)A_w]
%        =int_BZ |what(q)|^2 B_mu(q) d^3q/(2pi)^3,          (QV1)
% B_mu(q)=(1/4)int_BZ [E(k)-E(k-q)]^2
%       *[1-d(k).d(k-q)/(E(k)E(k-q))] d^3k/(2pi)^3.       (QV2)
% No spin or valley multiplier is missing: Tr_spin[P_-(k)P_+(k-q)]
% equals the displayed square bracket for the FULL4 operator. Indeed
% P_-(k)A_w(k,k-q)P_+(k-q)
%   =what(q)[E(k-q)-E(k)] P_-(k)P_+(k-q)/2.
% The q Fourier convention is what(q)=sum_x w_x exp(-iq.x).
% The identity extends as a finite vacuum-variance formula to l2 weights
% by Hilbert--Schmidt continuity, but l1 already suffices below and
% avoids any unbounded observable-domain inference.

% QV2. The variance is STRICTLY positive for every such nonzero weight.
% The zero set of B_mu(q) is exactly q in {0,pi}^3 modulo 2pi.
% At these eight transfers E(k-q)=E(k), so (QV2) vanishes. Conversely
% take q outside that set. E(k)^2-E(k-q)^2 is the nonzero analytic
% trigonometric polynomial sum_i sin(q_i) sin(2k_i-q_i). Thus E-E'
% is nonzero on a dense open set away from the finitely many possible
% nodes. If B_mu(q)=0, nonnegativity and continuity of its integrand
% force d(k)/E(k)=d(k-q)/E(k-q) throughout that punctured torus.
% For mu>0 the mass component then gives E=E', a contradiction. For
% mu=0, set k_i=0 and choose the other two components away from both
% sets of nodes. Equality of the i components forces sin(q_i)=0 for
% every i, again a contradiction. The omitted node sets have measure
% zero; no gap or mass-limit interchange is used.
% A nonzero l2 Fourier function cannot be supported on eight points.
% Therefore V_mu[w]>0 for every nonzero real l1 weight. In particular
% O_w does not annihilate the incoming sea. The extended uniform and
% staggered conserved charges have non-l2 weights; their zero variance
% is not a counterexample to this localizable-source assertion.

% QV3. A genuine angular test survives ANY normalized radial L1 readout.
% Fix one such REAL K and finite ell, with the actual POINT allocation
% rho_ell(y)=sum_x K_ell(y-x) :e_x:, K_ell(y)=ell^-3 K(y/ell).
% For a merely L1 kernel this is read as an operator-valued distribution
% through the following absolutely convergent smeared observables.
% Khat(0)=1 supplies a ball where Khat is nonzero, including for the
% fixed noninjective kernel used in the local-area preparation. Choose
% a nonzero real Schwartz test f whose Fourier transform is a smooth
% radial annulus profile times a real even spherical harmonic of degree
% two, supported inside that ball after scaling and inside |q|<pi.
% The angular average of f at every radius is zero. Its readout weight
% w_x=int f(y)K_ell(y-x)d^3y=(f*K_ell)(x) is real and l1, since
% sum_x |w_x|<=||K||_1 sup_y sum_x |f(x-y)|<infinity.
% The no-alias Fourier identity gives what(q)=fhat(q)Khat(ell q),
% so w is nonzero. Thus the actual angular source
%              int f(y)rho_ell(y)d^3y=O_w
% has strictly positive vacuum variance by QV1--QV2, despite zero mean.
% The cell-box allocation has the extra Fourier multiplier given by
% the cell indicator; it is nonzero in a ball about zero. Choose the
% same annular support inside that ball. The weight remains real l1
% and nonzero, so the same angular test and conclusion apply. No global
% inverse of K, injective filter or finite-support source is assumed.

% QV4. Exact scope of the quantum obstruction and allowed alternatives.
% Suppose the DECLARED unprojected continuum Gauss law is to hold as
% an operator constraint on a spherical physical sector containing the
% actual incoming sea, with every displayed constraint defined on that
% vector. Suppose its geometric fields are operator-valued radial N(r)
% and v(r) rhat, with a rotation-covariant ordering. The energy flux is
% then an operator-valued radial scalar: every zero-angular-average
% smearing f above annihilates it. This remains true for noncommuting
% radial coefficient operators and does not rely on positive curvature
% squares, a weak-g expansion or a choice between the momentum signs.
% The operator energy constraint would require O_w|Omega>=0. QV1--QV3
% instead give ||O_w|Omega>||^2=V_mu[w]>0. Hence NO such exact spherical
% operator sector can contain this incoming sea with the stated native
% source/readout. This is independent of the centre square-ordering
% obstruction and holds at every finite mu and g>0.
%
% The Hartree EXPECTATION constraint avoids this contradiction because
% <O_w>=0. It does not thereby furnish the full operator reduction,
% its angular fluctuation dynamics or a fixed first-G pair kernel.
% A radial angular projection would remove the tested observable by
% changing the source map; allowing nonspherical geometric operators
% could carry it but changes the stated geometric sector. Extra constraint
% terms or a weak expectation-only interpretation would need their own
% explicitly specified construction. QV7 below further excludes every
% normal state of the incoming representation, independently of scalar
% reference subtraction; merely omitting the sea vector from that normal
% physical sector does not evade it. No such repair is silently used.
% In particular this theorem does NOT exclude a specially constrained
% Hartree history with exactly radial expectation values, nor does it
% assume stationarity of any candidate history. The full unprojected
% common-transport existence and uniform time/clock/ell problem remains.

% QV5. Hydrodynamic suppression is distinct from an exact constraint.
% There is a quantitative all-mass check of the preceding distinction.
% For a fixed vector Q, as epsilon->0,
% B_mu(epsilon Q)/epsilon^4 -> C_mu(Q),
% C_mu(Q)=(1/8)int_BZ (Q.grad E)^2 |Q.grad(d/E)|^2 d^3k/(2pi)^3.
% This homogeneous quartic is strictly positive for Q!=0. For mu>0,
% Taylor expansion on the torus proves the formula. At mu=0 remove
% balls of radius 2epsilon|Q| about the eight nodes. On those balls
% the integral is O(epsilon^5|Q|^5); outside, the integrand divided by
% epsilon^4 is bounded by C_Q/dist(k,nodes)^2, which is integrable in
% three dimensions. Dominated convergence proves the SAME coefficient.
% Positivity follows near a valley by choosing k nonorthogonal to Q
% and, at mu=0, also not parallel to Q: both directional derivatives
% are nonzero on an open set. The same split gives
% B_mu(q)<=C_mu |q|^4 near zero.
% The coefficient is an explicit full-zone lattice integral. With
% s_i=sin k_i, c_i=cos k_i and A=sum_i Q_i s_i c_i,
% C_mu(Q)=(1/8)<A^2 sum_i Q_i^2 c_i^2/E^4-A^4/E^6>.
% Cubic symmetry gives C_mu(Q)=b4 sum_i Q_i^4+b22 sum_(i<j)Q_i^2 Q_j^2,
% b4=(1/8)<s_1^2 c_1^4(mu^2+s_2^2+s_3^2)/E^6>,
% b22=(1/8)<c_1^2 c_2^2[(s_1^2+s_2^2)E^2-6s_1^2 s_2^2]/E^6>.
% Node values are interpreted by their integrable limit. For large mu,
% <sin^2 k cos^4 k>=1/16 and <sin^2 k cos^2 k>=1/8 give
% C_mu(Q)=|Q|^4/(128mu^4)+O(|Q|^4 mu^-6), not a new species count.
%
% For a fixed real Schwartz angular f with fhat Khat nonzero as in QV3,
% set f_ell(y)=
% ell^-3 f(y/ell), W=K*f and w_(ell,x)=ell^-3 W(x/ell). At large ell
% the fixed compact support of fhat excludes aliases. Then
% ell^7 V_mu[w_ell] ->
%     int_R3 |fhat(Q)Khat(Q)|^2 C_mu(Q)d^3Q/(2pi)^3 >0.   (QV3)
% This is a fixed-mu, fixed-filter large-ell asymptotic, not a new
% numerical coefficient or a uniform collapse-state estimate. Finite
% averaging suppresses the variance but never makes the exact operator
% angular constraint true. In lattice units all coefficients depend on
% mu=ma and ell/a; no G occurs in this incoming-source variance. For
% the normalized density weight the physical variance is a^-8 times
% the displayed lattice variance. At every g=G/a^2>0 a radial Einstein
% flux still has identically zero angular smearing. A universal relative
% N_f^-1/2 fluctuation law is not supplied by the zero subtracted mean.
% The angular test need not live at the averaging scale. In the physical
% hierarchy 1<<ell<<L, take instead f_L(y)=L^-3 f(y/L). Its weight has
% Fourier transform fhat(Lq)Khat(ell q), and the same calculation gives
% L^7 V_mu[w_(L,ell)] -> int |fhat(Q)|^2 C_mu(Q)d^3Q/(2pi)^3>0
% for ANY ell/L->0 sequence. Continuity Khat(0)=1 and the compact
% Fourier support justify dominated convergence. Thus a macroscopic
% angular smearing also has nonzero variance at every finite L, while
% its leading suppression is independent of the averaging kernel in
% this hierarchy. This is a vacuum fluctuation result, not an emitted
% energy or thermodynamic entropy count and not exact finite-L equality.
%
% QV6. The obstruction also tests a specified physical Gaussian domain.
% At any slice keep the actual transported reference, but let the
% physical state be Gaussian with covariance C. Constant subtraction
% changes no variance, so for the same real angular weight
%       Var_C O_w=Tr[C A_w(I-C)A_w].
% If C,I-C>=eta Pi_B globally, then positivity twice gives
%       Var_C O_w>=eta^2 ||Pi_B A_w Pi_B||_HS^2.           (QV4)
% Choose a finite B with nonzero displayed compression. Such a B
% exists: for mu>0 a nonzero w_x gives an onsite mass block; at mu=0
% the absence of every nearest-neighbor block would imply
% w_(x+e_i)=-w_x for all x,i, impossible for a nonzero l1 sequence.
% Therefore any Gaussian state with these global margins on every
% finite set has positive angular-source variance and cannot satisfy
% the exact spherical operator constraint. This is an operator-form
% margin, not merely strict reduced covariance eigenvalues.
% The all-mass physical preparation in local-area UA14c has those
% margins. Their conditional persistence under an already existing
% bounded common-transport history is the separately paired AC44
% theorem, not a new existence assertion. No Gaussian formula is used
% for an arbitrary non-Gaussian reference or interacting state. This
% Gaussian argument alone leaves non-Gaussian states untested; the
% independently supplied normal-state theorem checked in QV7 closes
% that escape within the incoming representation, without a Wick
% formula for the physical non-Gaussian state's own covariance.

% SOURCES / INDEPENDENCE. Full binding GOAL_alt.md reread during
% this continuation. QV1--QV6 use the actual full4 sine symbol, native
% symmetrized energy vertex, common incoming sea and fixed point/cell
% readout; the projector trace, positivity set, Fourier test and scaling
% proofs are given above. Sections 14--15 and the published faithful
% and local-area source domains retain their separate scopes. This
% substantive save precedes any new original-RED or mechanics-side
% operator-angular continuation intake. No new counterpart angular result,
% numerical result/run or external X is used. The published core was
% unchanged at that independent save; the compact result now includes the
% original pair's actual-source comparison and normal-state extension.
% QV6 uses the already paired local-area UA14/AC44 order domain as
% explicitly attributed; its new double-positivity variance test is
% derived here before the original RED's new angular-operator intake.

% QV7. ACTUAL QR1--QR5 COMPARISON AND THE STRONGER NORMAL-STATE DOMAIN [R].
% Read the complete actual RED section 17 and SOURCES only after both
% substantive independent saves. Its interband/full4 formula agrees
% with QV1, and its separate small-transfer witnesses prove the needed
% positive variance at every finite mu. QV2 adds the exact eight-point
% zero set. The two independently derived quartic coefficients agree
% after Q_i=sin k_i cos k_i; the massless node split and macroscopic
% L^-7 limit also agree. QV1 is the infinite-parent form of the already
% paired section 5/QC2 variance identity; the new classification,
% angular constraint and scale conclusions are distinguished from that
% predecessor formula. The following stronger normal-state argument
% is independently supplied by RED QR2 and checked here, not recredited
% as a blind author derivation.

% Choose a REAL RADIAL Schwartz base test phi, with nonzero smooth
% Fourier support in a small ball where Khat(ell q) is nonzero, and
% define w=phi*K_ell sampled on the sites. QV1--QV2 give sigma^2>0
% for B=dGamma(A_w)-Tr(P0 A_w). For integer R put
%        B_R=tau_(R e_1)(B)-tau_(R e_2)(B).
% The actual continuum smearing is phi(y-R e_1)-phi(y-R e_2).
% Its angular average vanishes at EVERY radius because the two centers
% have equal distance from the origin. Thus any radial operator scalar
% flux would require the corresponding strong angular source constraint.
% No radial projection or new source allocation is introduced. With the
% cell-box allocation the base weight gains the cell Fourier factor;
% it remains nonzero near zero and trace class. The same translated
% RADIAL test has zero angular average regardless of that cell factor,
% so the normal-state proof extends to that allocation as well.

% The free sea is clustering: its bounded torus projector is integrable,
% so every matrix entry P0(x-y) tends to zero at large separation.
% Wick expansion then factorizes separated finite local CAR collections.
% The translated B's are even. For every fixed local CAR polynomial F,
% the three separated clusters centered at 0,R e_1,R e_2 therefore give
% <F Omega,B_R F Omega> -> 0,
% <F Omega,B_R^2 F Omega> -> 2 sigma^2 ||F Omega||^2.
% First use finite-site approximants to B, then its uniform operator-
% norm approximation from trace-class A_w. The norms ||B_R||<=2||B||
% are uniform. The vectors F Omega are dense; approximation and
% polarization extend both weak limits to the whole sea Hilbert space.
% Finite-rank approximation of a normal density matrix rho then gives
%       Var_rho B_R -> 2 sigma^2>0                       (QV5)
% for EVERY normal rho, with arbitrary coherences and higher correlations.
% Wick factorization was used only in the vacuum clustering step, not
% as an assumption that the physical state rho is Gaussian.

% Every strong constraint B_R-c_R=0 on the support of such a state
% would instead give Var_rho B_R=0 for every R. Thus NO nonzero normal
% physical sector in this representation satisfies all the exact
% unprojected spherical operator source constraints. The arbitrary real
% c_R emphasizes that changing the scalar normal-ordering subtraction
% cannot repair the result. This excludes normal non-Gaussian states
% as well as the incoming sea and QV6's Gaussian order-margin domain.
% All ordinary finite expected particle-plus-hole excitations are normal
% here: tails of their excitation count tend to zero outside increasing
% finite mode sets, whose vacuum-tail restrictions approximate the state
% in norm. More generally normality alone is the theorem's premise.
% A self-adjoint quantum Hamiltonian on this Hilbert space preserves
% normality; choosing a different interaction ordering within it does
% not create a nonzero strong-constraint sector. No stationarity or
% particular source evolution was assumed in this argument.

% The statement is infinite-parent and representation-specific. It is
% not an all-state theorem for an arbitrary finite box, an inequivalent
% nonnormal representation, or constraints asserted only after a new
% Hilbert-space compression. Compression can discard the angular source
% fluctuations; its native invariance, algebra and Hartree matching
% would have to be constructed. If the physical subspace is invariant
% under the uncompressed source constraints, vanishing compressed
% constraints implies the strong condition already excluded. These are
% named alternative premises, not adopted repairs of the binding map.

% QR4's conditional transported-representation extension also checks.
% Assume an ALREADY GIVEN finite-time bounded finite-range history,
% with integrable coefficient bounds, whose coefficients tend at
% spatial infinity to translation-invariant N_infty(t),v_infty(t).
% Require the translated one-body generators to converge strongly for
% a.e. t with an integrable norm bound. The limiting h_infty commutes
% with P0. Duhamel gives T_z^* U(T)T_z -> U_infty(T) strongly for
% both propagators and adjoints. For trace-class A_w this implies
% trace-norm convergence of the conjugated vertices. Pulling the actual
% normal ordering P_T=U P0 U^dag back to the initial representation
% therefore makes each distant angular observable converge in norm
% to a translate of the limiting centered vertex. Its sea variance
% is still sigma^2 because U_infty commutes with P0. The same two-center
% argument proves (QV5) for commonly transported normal initial states
% even if the one-body transport is not implementable in the original
% sea representation. No existence, clock bound or uniform ell/time
% estimate for such a history is inferred. Asymptotically flat bounded
% histories satisfy the named spatial limit if they are actually given.

% CONSEQUENCE AT THE BINDING SCOPE. Under duty 9's strong unprojected
% operator-functional reading, a radial geometry on the native incoming
% sea's normal Hilbert space cannot satisfy the declared energy Gauss
% constraints; the obstruction is angular and independent of either
% momentum sign, finite mu, rotation-covariant radial ordering and every g>0.
% It is stronger than a failure of a chosen square-root prescription.
% Hartree equations impose expectation values and are not contradicted
% by this theorem. A special unprojected Hartree history, its invariant
% coherence class and the required uniform time/clock/ell bounds remain
% to be established. The accessible quantum count, matched emitted
% Gauss energy and detailed-balance entropy also remain uncomputed;
% vacuum fluctuation suppression is none of those results. No external
% X or complete constructive quantum model is claimed by this comparison.

% FINAL ACTUAL ANGULAR-OPERATOR READ / SOURCES [R].
% Both complete current sources, declared_exchange_kernel.tex and
% redteam_declared_exchange_kernel.tex, including every SOURCES block,
% were read after the independent QV and QR saves and reciprocal comparison.
% The exact full4 interband trace, eight-point zero set, all-mass node
% estimates, quartic coefficient and Gaussian-margin test agree. The
% independent RED clustering/dense-vector proof was checked as the broader
% normal-state argument, with scalar-shift invariance and its conditional
% strong-star transported-representation transfer retained explicitly.
% The compact statement now uses differences of equal-radius translates,
% distinguishes a radial kernel from cell-box allocation, and requires
% rotation-covariant radial flux. No ordering change within the same
% normal representation removes the strong-constraint contradiction.
% The remaining Hartree mean-history, uniform clock/time/ell and physical
% emitted-energy/state-count problems are not claimed solved. A proposed
% compression, nonnormal representation or weak constraint needs its
% changed premises, invariant domain and actual Hartree matching.
% The RED's exact revised compact paragraph, projected-history wording,
% independence chronology and QR6 comparison were read as saved. No new
% external X, numerical check or complete quantum construction is inferred.

% QS1--QS4. PHYSICAL FLUCTUATIONS AND A SUPPLIED PAIR ON THE NATIVE SPIN
% CODE [R; independently saved continuation, compact source unchanged].
% The entropy pair's checked Pauli symmetry and common-mean code are the
% starting input, not recredited here. The following physical variance
% and quartic-compression calculations were derived before reading any
% new RED continuation on these corrections. They do not introduce an
% ordering/sector that evades QV7 or supply the missing Hartree seed.

% QS1. Fix an ACTUAL common bounded quadratic point-field history, if
% one is supplied. At one time, write the one-body factorization as
% H=K tensor C^2, P=P_T=p tensor I_2. Let w_1,...,w_n be orthonormal
% vectors in (I-p)K, W=sum_i|w_i><w_i|, and Pi the many-body code with
% exactly one particle in each orbital and zero holes above P. This
% describes the transported entropy code, without assuming a Fock
% implementer for the entire infinite-parent one-body U. The bounded
% CAR expectations below may equivalently be pulled back to the incoming
% representation. For any REAL native averaged energy/current test,
% let A=tilde A tensor I_2 be its self-adjoint trace-class vertex and
% O_A=dGamma(A)-Tr(P A). Put a_ij=<w_i,tilde A w_j>, q=sum_i a_ii.
% Then Pi O_A Pi=q Pi, but the physical variance need not be scalar.
% Let S_ij exchange logical spins i,j; it is Hermitian with S_ij^2=I.

% The exact variance operator on the code is
% Pi(O_A-q)^2 Pi
%  =[V0-L_-+L_+] Pi
%      +2 sum_(i<j)|a_ij|^2(I-S_ij),                    (QS1)
% V0=Tr_H[P A(I-P)A]=2||p tilde A(I-p)||_HS^2,
% L_-=sum_i||p tilde A w_i||^2,
% L_+=sum_i||(I-p-W)tilde A w_i||^2.
% In particular V0 is the ACTUAL transported-sea variance, retained
% despite its zero normal-ordered mean. It is not an excess variance
% with the sea contribution subtracted or clipped. All traces converge
% for the stated vertices; finite-support/finite-parent versions follow
% by the same finite-CAR algebra.

% Proof. Full reference Wick ordering gives the operator identity
% O_A^2=:dGamma(A)dGamma(A):_P
%       +:dGamma(A^2-2 A P A):_P+V0 I.
% Compress the quartic term onto one particle per orbital. Its diagonal
% terms give sum_(i!=j)a_ii a_jj, while its exchange terms give
% -sum_(i!=j)|a_ij|^2 S_ij. The contraction term gives
% sum_i <w_i,(tilde A^2-2 tilde A p tilde A)w_i>.
% Subtract q^2 and resolve I=p+W+(I-p-W) to obtain (QS1).
% This proof uses no Gaussian assumption on the logical spin state.
% A Gaussian state with the same one-body covariance does not in general
% have these fixed-occupation four-point functions.

% Since L_-<=V0/2 and I-S_ij>=0, every logical state rho obeys
% V0/2 <= Var_rho(O_A) <= V0+2n||A||^2.                (QS2)
% For the upper bound use I-S_ij<=2I and
% L_++4sum_(i<j)|a_ij|^2 <=2sum_i||tilde A w_i||^2.
% Thus even this source-blind code cannot remove the physical sea
% fluctuations. If the additional STATE/TEST premise q>=n e_A>0 holds,
% then sigma/q <=sqrt(V0+2n||A||^2)/(n e_A). This supplies a concrete
% n^-1/2 upper scaling only with its actual e_A, vertex norm, sea variance
% and geometric/averaging dependence controlled. Finite count alone
% supplies none of those uniform constants or a positive mass operator.
% With A=a^-1 Ahat, q=a^-1 qhat and V=a^-2 Vhat, the relative bound
% is dimensionless; g=G/a^2 and mu=ma enter the actual fields, orbitals
% and averaged test, not an invented universal coefficient of n^-1/2.

% QS2. A two-orbital witness distinguishes identical means physically.
% At n=2, choose a_12!=0. The code singlet has S_12=-1 and every
% triplet has S_12=+1. Their means q are identical, while
% Var_singlet(O_A)-Var_triplet(O_A)=4|a_12|^2.           (QS3)
% This is an actual physical-source statement whenever the chosen
% source vertex has that nonzero matrix element; it is not a claim
% that every test distinguishes every code. Such vertices and orbitals
% exist natively: any nonzero finite-site energy vertex has a nonzero
% positive-band compression for mu>=0, and that finite-rank compression
% cannot be a scalar on the infinite positive orbital space. Two
% orthogonal combinations of vectors with different quadratic forms
% therefore give a_12!=0. For the standard site vertex this nonvanishing
% also follows directly from its positive-band kernel
% (E(k)+E(k'))P_+(k)P_+(k')/2, which is not identically zero.
% This witness changes no mean-source branch. It exposes information in
% physical fluctuations that a Hartree expectation does not measure.

% QS3. Exact action of an ALREADY SPECIFIED first-g pair. Use lattice
% units and dimensionless vertices Ahat_alpha=a A_alpha. Suppose a
% separate valid construction has supplied a real symmetric c-number
% kernel K_alpha beta on a finite family of these actual vertices, or
% a sum convergent in operator norm with the stated trace contractions.
% At the given history
% time freeze that kernel and P, and declare FULL reference Wick order:
% a H_pair=(g/2)sum_(alpha,beta) K_alpha beta
%                  :Ohat_alpha Ohat_beta:_P,
% g=G/a^2. Symmetry makes the sum Hermitian; the displayed convention
% removes all reference contractions, not only the two mean sources.
% This is a test of a supplied pair, including the conditional kernels
% already derived above. It neither derives a universal first-g kernel
% nor drops the native sqrt(g) terms of a different expansion.

% Let u^alpha_ij=<w_i,tilde Ahat_alpha w_j>,
% V_ij=sum K_alpha beta u^alpha_ii u^beta_jj,
% J_ij=sum K_alpha beta u^alpha_ij u^beta_ji.
% Both V_ij=V_ji and J_ij=J_ji are real. Direct CAR compression gives
%       Pi(a H_pair)Pi
%          =g sum_(i<j)[V_ij I-J_ij S_ij].              (QS4)
% The two annihilators in an exchanged ordered pair contribute the
% minus sign; equal orbital indices vanish because the code contains
% only one particle there. Terms taking particles outside the occupied
% orbital set or creating holes are removed by Pi, not asserted absent
% from the full Hamiltonian. At n=2 the singlet/triplet matrix-element
% splitting is 2g J_12. It need not vanish when all Hartree sources
% agree. If the fixed kernel is the weak Newton prescription, its own
% sign must be retained in J_12; no positivity is assumed here.

% For a product of orbital spinors s_i, <S_ij>=|<s_i,s_j>|^2.
% The expectation of (QS4) then agrees with Hartree minus exchange
% for that Slater state, including cancellation of the Hartree self
% terms. For entangled logical inputs the exact swap expectation is
% needed; a Gaussian Wick reconstruction from their covariance can
% introduce false self-pairs and miss connected spin correlations.
% Thus the 2^n common-mean count is not automatically a degeneracy of
% a quantum interacting/boundary-charge operator. Global SU(2) symmetry
% of a pair also does not preserve every independently encoded spin.

% QS4. Scope of the compression. Equations (QS1)--(QS4) are exact
% physical expectations/matrix elements on a specified finite-count
% code at a supplied common quadratic history. Pi H_pair Pi is not
% automatically its effective dynamics: invariance, leakage control
% or a justified degenerate perturbation reduction would be separate
% conditions. It does not replace the uncompressed source operators,
% evade their strong angular no-go, or declare a new reduced sector.
% Nor does it construct a self-consistent unprojected Hartree seed,
% the geometry's operator functional, an emitted Gauss charge, or an
% accessible thermodynamic count. It computes a concrete physical
% fluctuation and the Fock/spin dependence for any genuinely supplied
% full-Wick pair at the explicit named scope. No code/run is required.
% SOURCES / INDEPENDENCE: the already paired native Pauli algebra and
% one-per-orbital code in declared_entropy.tex EC-A5; direct CAR/Wick
% multiplication, positive projections and swap algebra above. No new
% RED correction or counterpart pair source was read before this save.

% QS5. THE FIXED ORBITAL MIXTURE RETAINS ITS CLASSICAL SOURCE VARIANCE [R].
% Read RED entropy EC10 only after QS1--QS4 were saved. Its actual
% twirled-half-occupation, integer-N and fixed-weight encoding proof
% is checked in declared_entropy.tex. Write Aset for an N-orbital set
% in that preparation, avoiding confusion with the one-body test A.
% For its code mean q_Aset=sum_(j in Aset)<w_j,tilde A w_j>, the exact
% physical law of total variance gives
% Var_Enc(rho)(O_A)
%  =sum_Aset p_Aset Var_(V_Aset rho V_Aset^dagger)(O_A)
%     +sum_Aset p_Aset(q_Aset-qbar)^2,
% qbar=sum_Aset p_Aset q_Aset.
% It follows simply by expanding Tr rho O_A^2 and the squared mean;
% O_A need not preserve the orbital-set sectors. Each within-code term
% is exactly QS1 with the corresponding W_Aset. The second term is
% physical preparation variance, not a sea counterterm. The whole
% family has input-independent means, yet need not have a universal
% relative N^-1/2 fluctuation bound.

% An actual native initial-state example makes that last limit precise,
% without claiming a spherical or reached history. Set a=1 here. Choose
% a measurable normalized positive spinor on D=(-pi/2,pi/2]^3, take
% chi(k)=sqrt(8)u_+(k)1_D(k), and e=(chi+T_3 chi)/sqrt(2).
% The two supports are disjoint; e lies in K and in the positive band.
% Its translates e_x=T_(2x)e, x in Z^3, are exactly orthonormal, since
% the normalized Fourier transform over D gives delta_xy at even
% translations. Their common free mean energy Ebar is strictly positive
% at every finite mu>=0 and at most Emax=sqrt(mu^2+3). No spatial
% regularity beyond the actual square-summable vectors is presumed.

% Fix a positive normalized smooth radial averaging kernel with a
% bounded periodized lattice sum, as in the specified readouts already
% used here, and a finite ell. For a ball of radius R define the actual
% averaged mass vertex A_R={W_R,h0}/2,
% W_R(x)=int_(|y|<=R) K_ell(y-x)d^3y. Then 0<=W_R<=1,
% ||A_R||<=Emax, A_R is trace class, and A_R->h0 strongly as R->infinity.
% Choose any N distinct e_x. For sufficiently large R_N, their mean
% q_Aset for A_(R_N) is >=(1-epsilon)N Ebar. Keeping that ball fixed,
% choose N other even translates arbitrarily far away. Compactness of
% A_(R_N) and weak convergence of distant translates make their q_Bset
% satisfy |q_Bset|<=epsilon N Ebar. Both orbital sets have exact free
% mean energy N Ebar and no holes, and all their vectors remain mutually
% orthogonal. R_N may be taken much larger than the fixed ell and a.

% Mix the two one-per-orbital codes with fixed weights 1/2. Its twirled
% excess is (W_Aset+W_Bset) tensor I_2/4<=P_+/4, so it satisfies EC10's
% actual occupation and count premises. The extra mass variance is
% (q_Aset-q_Bset)^2/4 >=(1-2epsilon)^2 N^2 Ebar^2/4.
% For epsilon<=1/4 its positive mean is at most N Emax; hence its
% relative physical standard deviation is at least Ebar/(4 Emax),
% independently of N. This disproves a variance-suppression inference
% from finite count and those source-blind encoding premises alone.
% The example supplies no unprojected spherical initial charge, actual
% coupled trajectory or geometric scale bound, and thus is not a
% counterexample to a separately proved reached-state estimate with
% stronger localization/source conditions. For such states the actual
% within-code and orbital-set terms must instead be bounded. Restoring
% a multiplies both mass and its standard deviation by a^-1; their
% ratio is a function of mu and the specified preparation, with no
% invented g=G/a^2 dependence before a geometry is actually supplied.
% SOURCES: actual saved RED EC10, exact CAR source compression, elementary
% total variance, Fourier orthogonality and compact-vertex limits. No
% new pair kernel, interacting generator, numerical run or external X.

% ACTUAL QR7--QR10 COMPARISON / SOURCES [R].
% Read the complete saved RED QR7--QR10 after QS1--QS4's independent
% save and QS5's subsequent EC10 intake. The direct four-CAR square and
% the full-reference-Wick calculation agree, including V0-L_-+L_+,
% the logical swap, the factor 2 in the upper bound and the retained
% physical sea variance. The pair uses a supplied real symmetric kernel;
% restoring a H_pair and dimensionless vertices gives exactly QS4.
% The two-set example and the infinite interval-sampler domain were
% checked from their actual saved formulas, not from a count-only claim.
%
% QR9's additional conditional global-symmetry construction is sound.
% Since the native Pauli action commutes with the sea projection, its
% relative-Fock second quantization leaves that sea invariant. An ACTUAL
% self-adjoint dynamics commuting with this SU(2), with its invariant
% domain supplied, decomposes as I_(V_J) tensor U_J on each spin irrep
% and multiplicity factor. Bounded spin-scalar source tests have the
% same form. Fixing one normal multiplicity state therefore preserves
% the V_J logical state and its input-independent means. The symmetric
% n-spin part of the one-per-orbital preparation has J=n/2 and n+1
% logical states. Generic exchange swaps act on other multiplicities,
% so this is not preservation of the full 2^n Hartree code. Interacting
% excitation-count bounds, the actual transported reference, a genuine
% self-consistent seed, spatial recovery and emitted Gauss energy are
% still required. No invariant quantum generator, new strong-constraint
% sector or larger accessible object count is inferred from this
% representation-theoretic statement. No new external X or run is used.

% ACTUAL QR11 NORMALITY / RADIAL-ORDER COMPARISON / SOURCES [R].
% Read the complete actual QR11 after its independent save. Its finite
% excitation-count implication does not assume a Gaussian state or
% normality. On the first M particle/hole modes, the no-excitation
% projection for modes m+1,...,M fails with probability at most
% delta_m=sum_(j>m)omega(b_j^dagger b_j). Gentle compression costs
% 2sqrt(delta_m); partial tracing its vacuum block differs from the
% first-m density by a positive remainder of trace at most delta_m.
% The uniform finite-M bound therefore gives norm distance at most
% 2sqrt(delta_m)+delta_m from the first-m state with vacuum tail.
% Finite expected excitation count makes delta_m->0, and norm closure
% of normal states proves the claimed incoming-representation normality.
% Finite net charge is not sufficient. At a different transported
% polarization the statement is in that representation, and QV7's
% actual asymptotic-history transfer remains a separate premise.

% Also checked the ordering scope directly: for exactly spherical
% operator fields, any ordered K_ij=A(r)delta_ij+B(r)n_i n_j has
% quadratic scalar contractions involving only radial words A^2,AB,
% BA,B^2. Noncommuting radial factors, inverse clocks and radial
% derivatives do not change this angular tensor identity. A radial-only
% nonlocal kernel likewise adds no angular direction. Thus those
% declared orderings are covered by the equal-radius variance no-go;
% a changed ordering alone does not evade it on finite-excitation
% states. Fixed-axis counterterms, nonradial geometry or angular source
% projection would change the stated premises. This strengthens the
% domain reading of the existing strong-constraint result without
% excluding an unprojected Hartree MEAN history or constructing its
% full quantum dynamics. No new external X or scientific run is used.

% IR1--IR4. INCOMING-REPRESENTATION DOMAIN OF AN ISOLATED-MASS HISTORY [R].
% Independent substantive save before reading the original RED's new
% asymptotic-flow/representation continuation. This is a conditional
% theorem about the ACTUAL native common transport, not a prescribed
% field substituted for an unproved Gauss solution. The printed core
% was unchanged at that independent save; its current statement includes
% the original pair's actual-source comparison. No run is used.

% IR1. Fix the history and its asymptotic domain first. Use a=1,
% finite mu=ma>=0, fixed g=G/a^2>0 and a fixed finite time T>0 on the
% infinite native parent. Let the actual real point fields generate the
% bounded finite-range h(t)={N,h0}/2+sum_i{v_i,Q_i}/2, with an integrable
% uniform coefficient bound on [0,T]. For clarity the sufficient bounds
% below are uniform in t; integrable bounds give the same proof. Assume
% outside a fixed ball, for |alpha|<=2 and some eta>0,
% |partial^alpha[N(t,x)-1]|<=C_alpha r^(-eta-|alpha|),
% v_i(t,x)=-b(t)x_i r^(-3/2)+v_i^rem(t,x),
% |partial^alpha v_i^rem|<=C_alpha r^(-1/2-eta-|alpha|), r=|x|.
% Coefficients are restrictions of these smooth continuum fields to
% actual native sites; there is no angular source projection. Take real
% b in L1[0,T], bounded here, which is NOT zero almost everywhere.
% The declared reference is P_T=U(T,0)P0 U(T,0)^dagger.
%
% For a genuine positive isolated-mass radial Gauss history with
% M(r,t)->E_boundary(t)>0 and N->1, the leading coefficient is
% b(t)=sqrt(2g E_boundary(t)) on the inward branch. The displayed
% derivative remainder rates require corresponding actual source/field
% decay; finite excitation count or a spatial averaging hierarchy alone
% does not prove them. Either printed momentum sign has this leading
% flux relation. No equality of their energy gradients is presumed.
% Here E_boundary is a times the physical energy, r is physical radius/a
% and t is physical time/a. Thus all constants retain their mu,g and
% actual-history dependence, as well as the chosen fixed averaging map.

% IR2. A local matrix expansion, valid also at the massless point nodes.
% Put n=x/r, A_ij(n)=delta_ij-(3/2)n_i n_j, q_i(k)=sin(2k_i)/2,
% and define the bounded Hermitian off-band matrix
% B_n(k)=sum_(i,j) A_ij(n) q_i(k) partial_(k_j)P0(k).
% At mu=0 the derivative is taken off the eight nodes. Its singularity
% is O(|k-k_v|^-1), cancelled by q_i=O(|k-k_v|), so B_n is bounded
% and defines an ordinary torus Fourier multiplier. The isolated nodes
% add no distributional delta term. For each FIXED native displacement z,
% direct finite-range propagation and first spatial Taylor expansion give
% (P_T-P0)_(x,x+z)=r^(-3/2) [T_b(n)]_z+O_z(r^(-3/2-delta)),
% delta=min(eta,1)>0,
% T_b(n,k)=-int_0^T b(s) exp[-i(T-s)h0(k)] B_n(k)
%                              exp[i(T-s)h0(k)] ds.          (IR.1)
% No assertion of exponential decay of the massless P0 kernel enters.
%
% Here is the matrix-element proof of the remainder. Freeze the fields
% at x and use h_x^0(t)=N(t,x)h0+sum_i v_i(t,x)Q_i. All these constant
% symbols commute in time and preserve P0. In the finite-range Dyson
% series, every row/column of their propagators, and of the actual U,
% has uniformly finite exponential weighted l1 norm at fixed T. A path
% reaching outside the ball |y-x|<r/2 has exponentially small weight.
% Inside it expand each coefficient to first order about x. The linear
% insertions are (partial_j N){Y_j,h0}/2 and
% (partial_j v_i){Y_j,Q_i}/2, Y_j=y_j-x_j. In the covariance expansion
% their commutators with P0 are EXACTLY
% [{Y_j,h0}/2,P0]=(i/2){h0,partial_j P0}=0,
% [{Y_j,Q_i}/2,P0]=i q_i partial_j P0.
% These identities hold distributionally on the torus at mu=0 too.
% Constant Q_i phases commute with B_n and cancel in the conjugation.
% The lapse phase differs from the free one by O(r^-eta), which gives
% the corresponding extra power in (IR.1). The remaining spatial terms
% are bounded by second derivatives and products of first derivatives:
% O(r^-5/2)+O(r^-2-eta)+O(r^-3)+O(r^-2-2eta), along with the specified
% flow-gradient remainder O(r^-3/2-eta). All are within the stated
% O(r^-3/2-delta) bound. Polynomial path-length factors from Taylor's
% formula have finite weighted Dyson sums. In U P0 U^dagger one may
% bound each intervening P0 matrix entry by its operator norm; the two
% propagator legs localize its endpoints near x and x+z. This proves
% the fixed-z remainder without an unproved summable massless P0 tail.

% The leading multiplier is nonzero on an open set of n,k. Since B_n
% is off-band, each of its two blocks is multiplied in (IR.1) by
% exp(+/-2iE T) int_0^T b(s)exp(-/+2iE s)ds.
% The Fourier transform of nonzero compact-time L1 b is entire and
% cannot vanish on an open frequency interval. To see that B_n supplies
% such an interval at EVERY finite mu, take n=e1, k=(p,p,0),
% 0<p<pi/2. Then B_n=q(p)[-(1/2)partial_1+partial_2]P0 is nonzero:
% the differentiated Dirac vector is not parallel to
% (mu,sin p,sin p,0). Its E=sqrt(mu^2+2sin^2 p) varies on an interval.
% Continuity away from the nodes supplies a positive-measure set of
% momenta and normals with T_b nonzero. Consequently there is a finite
% set Z of native displacements for which
% F_Z(n)=sum_(z in Z)tr|[T_b(n)]_z|^2
% has strictly positive sphere integral c_Z=int_S2 F_Z dOmega>0.
% One finite set suffices; no full-z asymptotic interchange is needed.

% Squaring (IR.1) for those displacements and summing over lattice sites
% gives the explicit lower divergence
% sum_(r0<|x|<R)sum_(z in Z)tr|(P_T-P0)_(x,x+z)|^2
%                   =c_Z log R+O(1),                    (IR.2)
% with a sufficiently large fixed r0. Indeed T_b is linear in A, so
% F_Z is smooth in n (a quartic polynomial in its components); comparing
% each unit cube with its integral costs an integrable O(r^-4) error.
% The cross term with the remainder is O(r^-3-delta), also summable.
% Thus P_T-P0 is NOT Hilbert--Schmidt for any such nonzero history,
% however small a fixed g>0 may be. If b=sqrt(g) times a fixed time
% profile, c_Z is proportional to g; changing an actual family need
% not keep that profile or its constants fixed. No numerical c_Z is
% asserted, and no finite-parent extrapolation is silently supplied.

% IR3. The representation consequence has a direct CAR proof. For two
% projections P0 and P_T, the extended nonnegative trace identity is
% ||P_T-P0||_HS^2=Tr[(1-P0)P_T(1-P0)]
%                         +Tr[P0(1-P_T)P0].
% Hence at least one incoming particle/hole branch has infinite expected
% occupation in the transported quasifree SEA. This alone would not
% exclude a general non-Gaussian normal state with infinite mean number;
% the following determinant and normality arguments are essential.
%
% In that divergent branch fix any finite initial set of incoming modes.
% For a finite set F outside it, the sea expectation of the projection
% onto no excitations in F is det(I-C_F)<=exp[-Tr C_F], where C_F is
% the particle covariance P_T or the hole covariance I-P_T compressed
% to F. This is the finite CAR/Wick determinant. Increasing F exhausts
% an infinite-trace tail and makes the probability tend to zero. In
% contrast, in the incoming sea Fock representation the projections
% onto no branch excitations outside the first m modes increase strongly
% to I as m->infinity. Every NORMAL state has their probabilities tending
% to one, even if its expected number is infinite. The transported sea
% therefore has no normal realization in that incoming representation;
% in particular no unitary on that Fock space can implement this sea
% transport. The one-body U and its CAR-algebra automorphism still exist.
%
% The distinction also persists for every state normal in the transported
% sea representation. First take a vector which is a finite CAR polynomial
% applied to its sea. Approximate that polynomial in operator norm by
% one supported on finitely many INCOMING modes. The tail-vacuum
% projections are even and commute with this approximating polynomial;
% its expectation is bounded by the polynomial norm squared times the
% sea's vanishing tail probability. Approximation and then the dense
% CAR-polynomial vectors extend this to all normal transported states.
% Such a state cannot at the same time be normal in the incoming Fock
% representation. QR11's already checked finite-count normality proof
% applies in the ACTUAL transported polarization. Thus finite expected
% particles plus holes relative to P_T puts physical payload states in
% this transported normal class; Gaussianity of the payload is not
% assumed. It does not restore implementation in the incoming class.

% IR4. Precise relation to the remaining binding problem. This is an
% obstruction to a fixed incoming normal-Fock quantum realization of
% the stated INFINITE-parent isolated-mass history with the displayed
% asymptotic controls. It is not a strong-Gauss ordering argument and
% does not rely on the separate angular-variance no-go. Nor does it
% forbid an actual Hartree mean history merely because its sea state
% changes representation. Finite parents have unitary implementers;
% their actual parent limit and reference/count domain must be supplied.
% The binding transported reference already permits considering its own
% representation. An algebraic or transported-representation realization
% still needs its quantum time-evolution/domain specification; declaring
% that reference does not construct the full ordered Hamiltonian. The
% excluded EXTRA premise is one fixed incoming normal-Fock realization,
% not the transported-reference law itself.
% Finite count alone does not give the asymptotic derivative rates, and
% histories not in that domain are not excluded. A genuine unprojected
% finite-g spherical mean history, uniform positive clock/time/ell
% controls, physical receiver/Gauss-energy transfer and accessible
% counts remain unresolved. No stationary covariance, zero radial
% derivative, refitted state, HA coupling or source projection is imposed.
%
% SOURCES / INDEPENDENCE FOR IR1--IR4. Actual native endpoint h0,Q_i,
% common transported P0, bounded finite-range Dyson expansion, the
% explicit off-band first-gradient identity and finite CAR determinants.
% The normal-state extension uses the already paired QR11 finite-count
% result, not a new incoming-reference substitution. This substantive
% local-matrix/log-divergence/domain route was saved before reading the
% original RED's new asymptotic-flow representation derivation. Existing
% QS and QV statements keep their separate supplied-interaction and
% strong-constraint domains. No new external X, code or run is used.

% ACTUAL QR12--QR14 COMPARISON / SOURCES [R].
% Read the complete saved RED QR12--QR14 and SOURCES only after IR1--IR4
% and that independent route existed. Its uniform gradient limit
% r^(3/2) partial_j v_i -> -b(t) A_ij(n), flat coefficient limits and
% r^(3/2) Hessian N ->0 suffice for the same fixed-entry leading term.
% In the exact commutator only the finite hopping endpoints occur;
% affine lapse cancellation removes the potentially larger first lapse
% gradient. On offsets within r/2, its scaled endpoint remainder has a
% polynomial offset bound. Outside, bounded commutators and fixed-time
% exponential propagator tails dominate every power of r. This checks
% the massless local limit without asserting decay of P0's absolute
% kernel sum. The nonzero band-frequency factor, open angular cone and
% dyadic lattice count give its logarithmic LOWER bound. IR's explicit
% power-decay premises additionally give the displayed finite-diagonal
% coefficient plus O(1); that sharper remainder is not imported into
% the RED's weaker o-limit domain.
%
% Its separate generator-domain statement also checks. At a time with
% b(t)!=0 the same affine commutator gives [h(t),P0] not HS, hence
% ||(1-P0)h(t)P0||_HS=infinity. The usual incoming-sea normal-ordered
% quadratic would create on its vacuum a pair vector with precisely
% that squared norm. Scalar subtraction cannot change the pair term.
% This excludes that vacuum from the formal quadratic's domain; the
% finite-time state argument is separately needed to exclude an incoming
% unitary implementer, rather than assuming a domain omission suffices.
%
% RED's independent normality route uses the full Gaussian excitation
% covariance, including anomalous particle-hole blocks: on a finite
% mode set Var N=Tr gamma(1-gamma)+Tr kappa^dagger kappa<=2 Tr gamma.
% The generalized CAR contraction proves the inequality. Arbitrarily
% large tail means would make their vacuum probability tend to zero
% by Chebyshev, contrary to incoming normal-state tail tightness. This
% agrees with IR3's branchwise determinant route and retains the crucial
% Gaussian qualification at that step. IR3's subsequent dense-polynomial
% argument addresses all transported-normal states, not all states with
% infinite incoming mean. Both routes keep actual transported-reference
% means separate from the incoming sea pairs: those pairs are not added
% again to the gravitating source or called an emitted Gauss-energy loss.
% The actual positive-mass history and its exterior bounds remain supplied
% premises. Finite parents, changed polarizations and infinite-past
% scattering preparations need their own limits/domains; no full ordered
% quantum completion or Hartree existence theorem follows. No new X/run.

% ACTUAL QR15--QR16 COMPARISON / SOURCES [R].
% Read the saved RED's attributed IR comparison and its subsequent
% instantaneous-domain refinement. The latter is a further consequence,
% not a claimed independent pre-intake derivation. Finite-range matrix
% multiplication of IR2 at a tested time gives
% r^(3/2)[h(t),P_T(t)]_(x,x+z)
%   -> [h0,T_b(n,t)]_z-i b(t)[B_n]_z.
% The h-h0 coefficients tend to zero at every fixed shifted endpoint;
% the separate [h,P0] term is exactly the affine-gradient limit.
% This uses no time derivative of an asymptotic error. In the (+E,-E)
% block put I_t(omega)=int_0^t b(s)exp[-i omega(t-s)]ds. Then
% T_b,+-=-B_n,+- I_t and the commutator multiplier is
% -i B_n,+- [b(t)-i omega I_t(omega)], omega=2E(k).
% The bracket is entire in omega and has value b(t)>0 at zero.
% Hence it cannot vanish on the open band-energy interval used in IR2.
% The same finite-offset angular-cone argument proves [h(t),P_T(t)]
% non-Hilbert--Schmidt at such a time, at every finite mu including zero.
% For a projection P_T this gives
% ||[h,P_T]||_HS^2=2||(1-P_T)hP_T||_HS^2=infinity.
%
% With real finite-support 0<=chi_R<=1 tending pointwise to one,
% h_R=chi_R h chi_R tends bounded-strongly to h. Its centered finite
% CAR quadratic has exact transported-sea variance
% ||(1-P_T)h_R P_T||_HS^2. Strong convergence and the finite-rank
% characterization of the HS norm show this tends to infinity.
% Thus this cutoff sequence cannot converge on the transported vacuum
% to the usual global quadratic generator. This is stronger than merely
% observing its domain problem at the incoming polarization; it still
% excludes neither a self-adjoint domain omitting that vacuum nor a
% specified connection between changing representation spaces.
% It asserts no equality between the quadratic and the nonlinear
% boundary-mass functional. The actual transported-reference source,
% local observables, both sign conventions and all actual-history
% premises are retained. No new numerical run or external X is used.

% ACTUAL QR17 CONSTRUCTIVE DOMAIN COMPARISON / SOURCES [R].
% Read the complete further RED continuation as saved AFTER its QR16,
% and checked it directly; it is not recredited as a blind derivation.
% For the actual supplied bounded history, let R_t=P_T and Q_t=1-R_t.
% The excitation space K_t=Q_t H direct-sum conjugate(R_t H) has its
% ordinary antisymmetric Fock space F_t. The unitary
% U|_(Q_s H) direct-sum conjugate(U|_(R_s H)) maps K_s onto K_t;
% its exterior second quantization is an exact inter-fibre propagator.
% It preserves the vacuum, total excitation number and composition.
% Precisely, its CAR intertwining carries pi_s(A) to pi_t(alpha_U(A)),
% NOT to pi_t(A) at the same algebra argument. It therefore does not
% contradict the disjoint normal classes in IR3 or supply an implementer
% of alpha_U in one fixed incoming representation. Centered trace-class
% source vertices transform with their actual reference and retain all
% physical local moments, including the sea variance and non-Gaussian
% payload cumulants; their difference covariance is exactly U D_s U^dag.
%
% With norm-continuous bounded h(t), R_t is norm C1 and
% K_t=i[dot R_t,R_t]=Q_t h R_t+R_t h Q_t is bounded Hermitian.
% The solution i dot W=K_t W transports R_0 to R_t. Consequently
% u=W^dag U commutes with R_0 and obeys
% i dot u=tilde h_parallel u,
% tilde h_parallel=W^dag(h-K_t)W,
% h-K_t=Q_t h Q_t+R_t h R_t.
% In this chart the excitation generator is the second quantization of
% Q_0 tilde h_parallel Q_0 direct-sum
%                       -conjugate(R_0 tilde h_parallel R_0).
% Differentiating conjugate(u) proves the hole minus sign. Orthogonal
% pinching gives ||h_parallel||<=||h||. On each n-excitation sector
% the generator norm is at most n||h||; the direct sum is self-adjoint,
% the finite-excitation vectors are a common invariant core, and
% ||H_exc(t)Psi||<=||h(t)|| ||N_ex Psi|| on D(N_ex). The second-
% quantized propagator is strongly differentiable on that domain.
% Combining this chart evolution with the W fibre map recovers U,
% rather than replacing the physical point Hamiltonian by h_parallel.
%
% Thus a concrete moving-domain quadratic construction IS available
% under the supplied-history hypotheses. The connection is determined
% by the declared transported reference, not an additional field or
% refilling prescription. This resolves that particular conditional
% domain task, while the actual unprojected coupled history, uniform
% time/clock/ell bounds, full ordered interacting generator, boundary-
% mass/source matching and physical emission/counting tasks remain open.
% It holds for either printed mean-field sign if the named history
% exists. It is no assertion that a nonlinear geometry operator has
% already supplied a one-body U. No new code/run or external X is used.

% BG1--BG4. THE FULL CURRENT-DEPENDENT MASS RESPONSE AND FINITE COUNT [R;
% independent substantive continuation, compact core unchanged].
% Derived before reading the original RED's new QR18 source-domain
% continuation. This uses the now paired IR/QR representation results
% as stated inputs. It tests a frozen response observable on an ACTUAL
% supplied background; changing a test state below does not assert that
% it continues to source that same geometry. No new history is inferred.

% BG1. Fix the infinite parent, finite mu>=0, g>0 and one actual bounded
% common-transport history satisfying IR1 through a time t>0 with b(t)>0.
% Fix a real nonnegative normalized radial Schwartz averaging kernel K,
% with point or cell-box allocation. A noninjective such K is allowed;
% no angular source projection is made. Let W^dag be its actual adjoint
% mapping continuum test fields to native site coefficients, and let
% B_e(f)={f,h0}/2, B_q(z)=sum_i{z_i,Q_i}/2. Assume the clock is positive
% and bounded away from zero at this time. Choose sigma=+1 for the
% compatible momentum sign and sigma=-1 for the literal sign. On the
% corresponding regular radial source branch, the exact linearized
% enclosed-mass variation already derived above is
% (N delta M)'=4pi r^2 N(delta rho+sigma u delta J), u=v/N.
% Its actual native response vertex for radius R is therefore
% A_R=N(R)^(-1){ B_e[W^dag(1_(|y|<R)N)]
%                     +sigma B_q[W^dag(1_(|y|<R)v)] }.       (BG.1)
% This retains the current and actual source-average placement. It is
% not the energy-only A_R of QV, nor the nonlinear mass operator itself.
% It is defined by the fixed branch's variation on radial sources; its
% finite CAR observable can also be tested in other physical states,
% without asserting that they obey the same spherical constraints.

% The coefficients in (BG.1) are summable at each finite R: the fixed
% Schwartz kernel has a bounded lattice periodization, and its continuum
% tests have finite-volume support and bounded amplitudes. Thus A_R is
% self-adjoint trace class. For sufficiently large R, its norm is
% uniformly bounded because N(R)->1 and W^dag is bounded on bounded
% fields. Dominated convergence gives bounded strong convergence to
% A_infinity=B_e(W^dag N)+sigma B_q(W^dag v).             (BG.2)
% The same conclusion holds with the stated cell allocation. The fixed
% kernel and cell box have finite spatial moments. Their convolution
% preserves the IR leading flow gradient and the vanishing scaled lapse
% Hessian: one splits |y-x|<|x|/2, uses the spatial Taylor bounds there,
% and uses the Schwartz tail outside. Consequently no equality between
% point and adjoint coefficients is assumed, but their asymptotic
% commutator coefficients used next DO agree. These statements are at
% fixed ell; uniform estimates as ell varies require their own bounds.

% BG2. Put P=P_T(t), and let T_b(n,t) be IR2's leading transported
% projector symbol. Finite-range multiplication gives, for each fixed z,
% r^(3/2)[A_infinity,P]_(x,x+z)
%    -> [h0,T_b(n,t)]_z-i sigma b(t)[B_n]_z.             (BG.3)
% The lapse first-gradient commutator with P0 vanishes exactly, as in
% IR2. In the (+E,-E) block the multiplier is
% -i B_n,+- F_sigma(omega),
% F_sigma(omega)=sigma b(t)-i omega int_0^t b(s)
%                                      exp[-i omega(t-s)]ds,
% omega=2E(k). It is entire with F_sigma(0)=sigma b(t)!=0. The already
% proved open-energy/cone argument therefore gives
% ||(1-P)A_infinity P||_HS=infinity                     (BG.4)
% for BOTH printed signs and every finite mu, including zero. This
% is a genuine current-dependent response statement. It neither drops
% the current nor replaces the mass derivative by the matter energy.
% Bounded strong convergence and HS lower semicontinuity now give
% V_R=Tr[P A_R(1-P)A_R] -> infinity as R->infinity.      (BG.5)
% Each finite-radius observable is legitimate. The divergence concerns
% its infinite-parent boundary limit at the stated actual background.

% BG3. Finite actual excitation count does not even ensure a finite
% boundary-response MEAN on the entire finite-excitation core. Pull
% (BG.1) back by the actual U(t,0), put Q0=1-P0 and
% Z_R=Q0 U^dag A_R U P0. Each Z_R is HS, while
% sup_R||Z_R||_HS=infinity by BG2. On the REAL Hilbert space of HS
% maps T:P0 H->Q0 H the continuous functionals
% f_R(T)=Re Tr(T^dag Z_R) have these unbounded norms. The uniform
% boundedness principle supplies ONE fixed T, normalized to ||T||_HS=1,
% for which sup_R|f_R(T)|=infinity. One may fix the integer radii first;
% the resulting test state is not refitted with radius or time.

% Let |Phi_T> be the normalized incoming one-particle/one-hole vector
% with pair amplitude T. For any 0<epsilon<1 define
% |Psi_0>=sqrt(1-epsilon^2)|Omega_0>+epsilon |Phi_T>.
% It is a physical neutral CAR vector with exactly
% <N_particles>=<N_holes>=epsilon^2, supported only in excitation
% sectors 0 and 2. Its initial positive free excitation energy is at
% most 2 Emax epsilon^2 and is strictly positive. Hence it lies in the
% common finite-excitation core of QR17, with arbitrarily small finite
% count and finite incoming free energy. Transport it by QR17's actual
% inter-fibre U; the same count and number-sector support are preserved.

% For O_R=dGamma(A_R)-Tr(P A_R), the transported mean is
% <O_R>=2epsilon sqrt(1-epsilon^2) f_R(T)
%       +epsilon^2 <Phi_T|:dGamma(U^dag A_R U):_(P0)|Phi_T>.
% The second term is uniformly bounded in magnitude by
% 2epsilon^2 sup_R||A_R||: within the one-particle/one-hole sector only
% the block-diagonal excitation quadratic contributes to its mean.
% The first term is unbounded along the chosen radius sequence.
% Therefore this ONE finite-excitation physical state has no finite
% boundary-response mean from the natural radius exhaustion. No
% Gaussian approximation or prior normality assumption was used.
% Its covariance is automatically CAR admissible. Its mean sources
% need not be radial and need not produce the supplied background;
% this is a domain counterexample, NOT an actual coupled-history
% counterexample or a no-go for a selected invariant spherical subset.

% BG4. Allowed restricted domains and physical variance remain distinct.
% If the actual excess D_0 is trace class, then D_t=U D_0 U^dag is
% trace class and bounded strong convergence does give the finite mean
% Tr(A_R D_t)->Tr(A_infinity D_t). Positive added-particle states with
% finite expected count have D_t>=0 trace class and satisfy this
% sufficient condition. Finite particle-plus-hole count alone gives
% only an HS excess in general, as the preceding coherent example shows.
% This is a concrete allowed mean-source domain, not an assertion of
% its exact spherical invariance or actual Gauss-history existence.

% Even that mean condition does not remove the physical sea variance.
% For any transported-normal state with ZERO holes and finite expected
% particle count N, write Gamma for its particle one-body covariance.
% Decompose the centered bilinear into its particle block and its
% pair-creation/annihilation blocks. On the zero-hole sector, the two
% output sectors are orthogonal. Direct CAR contraction gives
% Var(O_R)=Var_particle[dGamma((1-P)A_R(1-P))]
%                +V_R-Tr[Gamma(1-P)A_R P A_R(1-P)]
%          >=V_R-N||A_R||^2.                            (BG.6)
% The identity holds when the particle-block second moment is finite;
% if it is infinite the lower bound is automatic by truncation. It
% assumes no Gaussian particle state, no fixed-orbital code and no
% subtraction of variance. Thus BG5 makes this boundary-response
% variance diverge for every such zero-hole finite-count payload.
% A count-only relative fluctuation assertion is not supplied by the
% finite source mean or by the moving-Fock construction.

% The quantum response vertex is not asserted equal to the full
% nonlinear ordered mass function. A successful full construction
% would need its actual source/geometry response, invariant domain and
% boundary-generator relation; QR17 alone supplies none of those.
% The radius limit is after fixed a,mu,g,ell,t; finite-parent and joint
% hydrodynamic limits need separate control. In lattice units the
% displayed masses are a times physical masses, g=G/a^2 and mu=ma;
% b(t)=sqrt(2g M_infinity(t)) supplies the actual coupling dependence.
% No universal coefficient, g->0 interchange, emission rate or entropy
% count is inferred. SOURCES / INDEPENDENCE: the already paired radial
% integrating-factor response and IR/QR local-matrix theorem, exact
% finite-CAR particle/hole decomposition, bounded-strong convergence,
% HS lower semicontinuity and uniform boundedness. This substantive
% save precedes new original-RED source-domain intake. No run or X.

% ACTUAL QR18--QR19 SOURCE / BRANCH COMPARISON / SOURCES [R].
% Read both complete saved RED continuations after the independent BG
% save; their new implications are attributed to those actual sources.
% QR18 fixes a real normalized radial C1 kernel whose kernel AND first
% derivatives have bounded (1+|y|)^p weighted supremum, p>4, at the
% actual fixed ell. A cell box is included by its actual convolution.
% Thus C1 regularity alone is not claimed to give the required tail.
% Its initial physical excess obeys w_p D_0,D_0 w_p in S1,
% w_p=(1+|X|)^p. Finite-range conjugated propagation directly gives
% the displayed weighted trace-norm bound. The absolute diagonal sum
% of a trace-class matrix, applied after each native fixed shift,
% bounds the weighted energy/current arrays. The inequality
% 1+|y|<=(1+|x|)(1+|y-x|) then proves the actual convolved source
% and first-derivative bounds O(r^-p), uniformly on the supplied
% bounded interval. No spherical projection or spectral refill enters.
%
% Independently checking the actual radial rows, Y=sqrt(1+M) gives
% |Y'|<=2pi r^2|rho|+2pi sqrt(2g)r^(3/2)|J|. A finite baseline and
% nonnegative regular M therefore bound M and produce its finite limit.
% The separately required M_infinity>=E_*>0 makes the exterior clock
% integration regular. Its estimates N-1=O(r^(5/2-p)),
% N',N''=O(r^(3/2-p)), and the flow remainder estimates
% v_rem=O(r^(5/2-p)), partial v_rem,partial^2 v_rem=O(r^(3/2-p))
% do imply every IR1 bound with eta=p-4. In particular M'' retains
% the actual rho'=O(r^-p) term; no unjustified extra derivative decay
% is used. The second flow derivative sets this conservative exponent.
% Both printed signs have the same absolute estimates. The result is
% a sufficient source-to-exterior implication CONDITIONAL on the
% actual unprojected radial history, positive boundary mass and fixed
% scale. It supplies no interior-clock or uniform allowed-ell bound,
% no history existence and no source-only route to exact radiality.
%
% This domain makes the BG distinction concrete: its D_t is trace
% class, so the actual response mean has the finite limit in BG4,
% while the sea and every zero-hole finite-count payload retain the
% divergent boundary-response variance in BG2/BG4. The coherent HS
% but non-S1 domain counterexample in BG3 is not asserted to satisfy
% QR18's stronger source condition or to be an actual radial seed.
%
% QR19 subtracts the two already paired fixed-entry IR expansions for
% histories with the SAME P0 and asymptotic h0. If b_2-b_1 is nonzero
% on a set of positive time measure, its entire time transform gives
% a nonzero leading symbol and a finite-displacement logarithmic HS
% divergence for P_(T,2)-P_(T,1). This uses the whole tail history,
% not an unsupported inference from final boundary masses alone.
% The same determinant/dense-CAR normal-class proof applies to these
% two polarizations. Finite expected particle-plus-hole count relative
% to either branch cannot make their normal classes overlap. This is
% consistent with QR17's inter-fibre intertwining at DIFFERENT algebra
% arguments and with its conditional same-history spin code.
%
% Finally, every native CAR observable in the direct-sum representation
% is block diagonal, so a vector's relative branch phase gives exactly
% the same CAR state as the mixture of its two diagonal weights. This
% does not exclude an explicitly enlarged observable algebra, but none
% is supplied by an abstract Hilbert-space identification. It neither
% assigns entropy to boundary-mass branches nor proves a full quantum
% geometry or emitted charge law. No finite-parent limit, new X or run
% follows. All actual QR18--QR19 proof and SOURCES lines were read.

% BG5. ACTUAL QR20 ALL-FINITE-COUNT EXTENSION / SOURCES [R].
% Read the complete further RED result, derived there AFTER actual BG
% intake, and checked its CAR proof directly. This provenance is not
% an independent pre-intake derivation of that strengthening. In a pure
% sea fibre P, set Q=1-P and, for A=A^dag in S1, put
% Z=QAP, V=||Z||_HS^2, ||A||<=H. Let an arbitrary physical normal
% state have n_p=<N_p>, n_h=<N_h>, n=n_p+n_h<infinity and
% m_*=min(n_p,n_h). Write its centered bilinear as D_A+B^dag+B,
% with B=sum_j s_j b_j a_j in a singular basis of Z. Direct CAR gives
% [B,B^dag]=V-dGamma(ZZ^dag)-dGamma(conjugate(Z^dag Z)).
% Its mean is at least V-H^2 n. Vector Cauchy--Schwarz gives
% <B^dag B><=V sum_j<n_aj n_bj><=V m_*, and hence
% <BB^dag><=V(m_*+1). Applying the same argument to rho^(1/2)
% proves these statements for mixed and non-Gaussian states as well.
%
% For Y=i(B^dag-B), Var Y<=V C_m with
% C_m=(sqrt(m_*)+sqrt(m_*+1))^2. The remaining diagonal-block
% commutator is i(B_S^dag+B_S), S=A_++ Z-Z A_--. Since
% ||S||_HS<=2H sqrt(V), its mean magnitude is at most
% 4H sqrt(V m_*). The bounded-observable uncertainty inequality
% therefore proves, for V>0,
% Var(O_A)>=[V-H^2 n-2H sqrt(V m_*)]_+^2/(V C_m).     (BG.7)
% Finite-radius A_R is trace class, so both observables are bounded
% CAR operators and no hidden common unbounded domain is required.
%
% With the actual BG1 vertices and uniform H, V_R->infinity gives
% liminf Var(O_R)/V_R>=1/C_m>0 for EVERY fixed finite expected
% particle-plus-hole count state in the transported fibre. Arbitrary
% pair coherence cannot remove the divergence with finite count.
% BG4's zero-hole bound is sharper at finite R; BG7 extends its
% limiting conclusion to both species and to number fluctuations.
% A finite nonzero limiting response mean then has divergent relative
% standard deviation. This is not a universal finite-object coefficient
% and supplies no N_f^(-1/2) law. A radius-dependent state sequence or
% infinite actual excitation count requires a separate estimate.
%
% The tested observable is still the FULL linearized mass response
% on the fixed actual branch, not a proved nonlinear quantum charge.
% A physical test state need not retain that branch's radial sources.
% Thus this is a concrete constraint on a proposed common boundary-
% observable domain; it does not replace the missing actual seed,
% interaction ordering or boundary-flux relation. Sources: actual
% QR20/SOURCES, direct singular-basis CAR and the displayed variance
% inequality. No state projection, reference refill, run or X is used.

% BG6. ACTUAL QR21 PHYSICAL RADIAL-TANGENT COMPARISON / SOURCES [R].
% Read the complete new RED QR21 after the independent BG and UA17
% saves; this construction is attributed to that actual source. Its
% all-mass initial pinching preparation has C_*,1-C_*>=F/2 with a
% strictly positive polynomial lower bound on F, and exactly radial
% positive averaged rho_*, J_*=0 for the fixed noninjective point
% kernel. It is not an asserted reached history. Direct full4 native
% tracing of delta D=sum_i{g_i,Q_i}/2 gives delta e=0 and
% delta j_i(x)=[2g_i(x)+g_i(x+2e_i)+g_i(x-2e_i)]/8.
% Cross-current Gram terms vanish by their distinct two-step supports.
% The Fourier multiplier is therefore cos^2(q_i)/2, as claimed.
% The trace of P0 at every nonzero separation is zero in spin space;
% hence both initial particle and hole count changes vanish separately.
%
% Taking ghat_i=2i q_i fhat/[Khat(ell q)cos^2(q_i)] on the stated
% smooth small annulus produces exactly delta J=grad f and delta rho=0
% under the ACTUAL readout, with no reciprocal image in its pass band.
% The Schwartz g_i and polynomial F margin bound the weighted relative
% perturbation, so a fixed two-sided sufficiently small amplitude is
% CAR admissible. The displayed relative-mass contraction has a bounded
% integral coefficient, O(r^2) at the centre, and a finite exterior
% limit; it proves a positive regular initial branch for either sign.
% The clock integral is regular at both endpoints. This establishes
% actual physical radial variations, rather than assuming arbitrary
% independent continuum source variations are realizable.
%
% Independently checking the point pairing, the boundary derivative is
% sigma int v.grad f while Tr(h delta D) pairs the sampled v with the
% deconvolved native current. Their difference is exactly QR21's
% integral of [vhat_per(-q)/Khat(ell q)-sigma vhat(-q)].i q fhat.
% The r^-1/2 radial flow has vhat(q)=i b c0 q|q|^-7/2+O(1), c0>0.
% Its remainder r^-11/2 is integrable. The alias bound can be made
% explicit by first multiplying v by a smooth radial cutoff at radius
% R: the fourth-derivative cutoff errors are O(R^-3/2) in L1. Thus
% the nonzero reciprocal images have a uniform summable |G|^-4 bound;
% removing the cutoff proves the stated distributional Poisson formula
% against the annulus test. No continuum replacement of the lattice
% source pairing is made. The coherent leading radial integrand is
% b c0 m2_ell |q|^(1/2)/6, whereas bounded aliases are O(|q|).
% The literal leading term is 2b c0|q|^-3/2. A sufficiently small
% fixed nonnegative radial Fourier annulus therefore detects both.
%
% This is a finite-ell nonzero AFFINE boundary-gradient mismatch on a
% genuine physical exact-radial initial-data chart. The variation is
% not asserted isospectral, and a covariance Hamiltonian gradient can
% have commuting ambiguities on a chosen orbit. No universal quantum
% generator or actual-history no-go follows from this affine statement
% alone. A special invariant radial subset, a controlled vanishing
% mismatch in a scale limit, and the actual coupled history remain
% separate questions. The point coupling is not replaced by its adjoint.
% Sources: actual QR21/SOURCES, explicit native path traces, weighted
% CAR form bounds, the radial contraction and Fourier/alias check above.
% No additional external agreement, numerical run or new model is used.

% BG7. ACTUAL QR22 OFF-BAND / COVARIANCE-VELOCITY COMPARISON [R].
% Read the complete new original-RED QR22 and SOURCES after its
% substantive save and the author's UA18 save. The following checks
% and D-coordinate consequence are AFTER that intake, not an alleged
% independent pre-intake derivation. The published compact is unchanged.

% Retain exactly QR21's fixed noninjective Schwartz point readout and
% initial source profile, but fix mu>0 and 0<F<=1/4. The covariance
% C_*=Q0 F Q0+P0(1-F)P0 has free-band spectra in [0,1/4] and [3/4,1].
% Its global form margins C_*,1-C_*>=F/2 follow from
% F<=2(P0 F P0+Q0 F Q0), together with 1-F>=F.
% No pure-state or incoming-time evolution premise is inserted.

% For B_q(g)=sum_i{g_i,Q_i}/2 and
% delta D=Q0 B_q(g)P0+P0 B_q(g)Q0, direct endpoint multiplication
% independently verifies QR22.1. With p=k+q and
% v_i(p,k)=[Q_i(p)+Q_i(k)]/2,
% delta ehat(q)=0,
% delta jhat_i(q)=sum_j Xi_ij(q)ghat_j(q),
% Xi_ij(q)=2< v_i(k+q,k)v_j(k+q,k)
%                         [1-dhat(k+q).dhat(k)]>.
% Here the brackets mean the normalized FULL Brillouin-zone integral;
% the prefactor 2 is the sum of the two full4 off-band traces, not a
% species multiplicity. In the energy trace the two factors are
% (E_p-E_k)/2 and its negative with identical projector overlap, so
% they cancel BEFORE integration. The current vertex is spin scalar.

% The strict matrix positivity uses the mass premise. At midpoint l,
% v_i=cos(q_i)Q_i(l). For nonzero small q and mu>0, equality of dhat
% at the two endpoints implies equality of their energies and of all
% sine components. At least one cos(l_i)=0 is then necessary, a null
% set. Thus the weight is positive almost everywhere. For any nonzero
% complex vector z, the independent trigonometric functions Q_i(l)
% make |sum_i z_i cos(q_i)Q_i(l)|^2 nonzero on a set of positive
% measure. Xi is positive definite. Its real even smooth inverse is
% bounded on a FIXED small closed annulus away from q=0. Hence
% ghat=Xi^(-1)i q fhat/Khat(ell q) defines real Schwartz coefficients,
% and gives exactly delta rho=0, delta J=grad f with no aliases.
% There is no inverse asserted at q=0 or at zero mass.

% At fixed mu>0, analyticity of P0 in a narrow complex torus strip
% gives an exponentially decaying lattice kernel and boundedness on
% every polynomially weighted l2 space. Factor, for example,
% F^(-1/2)Q0 B_q(g)P0 F^(-1/2)
% =(F^(-1/2)Q0 F^(1/2))
%  (F^(-1/2)B_q(g)F^(-1/2))
%  (F^(1/2)P0 F^(-1/2)).
% The actual two-sided polynomial bounds on F and Schwartz g make
% every factor bounded. This proves the two-sided CAR perturbation
% domain, not just formal positivity of a Fourier symbol. Trace class
% and the off-band condition give zero change of each expected
% incoming particle and hole count. The actual radial affine path
% and both-sign regular Gauss contraction are those checked in BG6.

% Put X=Q0 delta D P0. Its trace-class Sylvester integral
% I_X=int_0^infinity exp(t C_+)X exp(-t C_-)dt converges with
% ||I_X||_1<=2||X||_1 and satisfies I_X C_--C_+ I_X=X.
% The Hermitian trace-class A with blocks A_+-=i I_X,
% A_-+=-i I_X^dag therefore obeys delta D=-i[A,C_*].
% This is an actual bounded-generator physical covariance-orbit
% tangent. The orbit path is not asserted to stay radial beyond
% first order; the distinct affine path establishes the exact radial
% differential. Neither path is substituted for a coupled history.

% Since the source changes coincide exactly with BG6, its native
% point/adjoint Fourier pairing gives
% Tr[(A_infinity^M-h)delta D]!=0 for BOTH signs, with
% A_infinity^M=h0+sigma B_q(W^dag v), h=h0+B_q(v).
% If a bounded self-adjoint H_E represented the boundary differential
% on all these physical exact-radial initial variations and also
% satisfied [H_E,C_*]=[h,C_*], then K=H_E-h would commute with C_*.
% The separated spectral projection 1_([1/2,1])(C_*)=P0 would commute
% with K, forcing Tr(K delta D)=0, a contradiction. This explicitly
% closes the commuting-gradient ambiguity left by the affine test.

% There is also a precise D-coordinate consequence of the same data.
% D_*=Q0 F Q0-P0 F P0 has strictly positive quadratic form on Q0
% and strictly negative form on P0, because F(x)>0 everywhere.
% It has no kernel; zero can remain an accumulation point. Therefore
% 1_((0,infinity))(D_*)=Q0 and 1_((-infinity,0))(D_*)=P0.
% Any bounded K commuting with D_* also commutes with these spectral
% projections. The SAME contradiction excludes a bounded H_E that
% both represents ALL the stated radial affine differentials and
% obeys [H_E,D_*]=[h,D_*]. No gap in D_* is needed for this spectral
% implication, and no bounded generator for delta D on the D orbit
% is inferred. A Hamiltonian defined only on a narrower radial orbit,
% another reduced bracket or an unbounded domain requires a separate
% argument; it is not excluded by demanding extra affine derivatives.

% The result concerns a nonempty fixed massive initial-state domain,
% with finite counts, exact point coupling and the actual both-sign
% Gauss variations. It neither proves that the native velocity is
% radial nor excludes every special invariant subset or finite-g
% history. SOURCES: actual saved QR22/SOURCES, the direct native spin
% traces, analytic massive projector bounds, Sylvester integral and
% spectral projections checked above, plus the already paired BG6
% Fourier/alias test. No new external X, run or changed dynamics.

% RR1--RR7. CONSTRUCTIVE REDUCED SPECTRAL-REFERENCE FORM [R, independent
% substantive derivation; original-pair comparison pending].
% The complete revised GOAL_alt.md and both REVIEWA_REVIEW files
% were read, including the owner overrides, before this save. No new
% counterpart revised-reference derivation was read. These formulas
% use the now-binding coherent sign, spherical source projection,
% adjoint site coupling and continued occupied spectral reference.
% They do not impose unprojected source equations or literal sea
% transport. The preceding strong-operator result is not continued.
%
% RR1. Fix a finite native parent and real source coordinates I=(x,alpha),
% with the actual four-component e_x and two-step j_(x,i) as vertices
% B_I. Let W map their cell charges to the smooth ANGULARLY AVERAGED
% rho(r),J(r), using the declared scalar and radial-vector kernels.
% Include their volume factors in W. Its genuine adjoint determines
% phi=W^dag(N,v), h=sum_I phi_I B_I. Thus no continuum sampling map
% or extra angular multiplicity is inserted in the matter operator.
% On a regular inward positive-mass branch, u=v/N<0,
% M=ru^2/(2G), M'=4pi r^2(rho+uJ),
% (log N)'=-4pi Gr J/u, M(0)=0 and N(infinity)=1.
% The branch may first be defined at a finite outer boundary, with
% the exterior limit justified separately by its actual source tails.
% Its linear identity (N delta M)'=4pi r^2 N(delta rho+u delta J)
% gives delta E=int dV(N delta rho+v delta J).
% This is the coherent sign only; the literal-sign results above are
% comparison controls and do not enter the reduced generator here.
%
% Let P(h) be the occupied spectral subspace CONTINUED on the chosen
% configuration chart. It has fixed occupations n_a=0,1, [P,h]=0 and
% P^2=P; an energy passing through zero never changes n_a. Locally
% require separation of oppositely occupied spectral clusters.
% A gap at the number zero is not required. An opposite-occupation
% degeneracy needs an explicitly resolved continuation; no divergent
% divided difference is used there. At finite parent, differentiating
% P^2=P makes delta P off diagonal in the P/(1-P) decomposition,
% whereas h is block diagonal. Hence Tr(h delta P)=0 exactly.
% With q_I=Tr[B_I(C-P(h))], the solved boundary functional obeys
%                 delta E[C]=Tr(h[C] delta C).                 (RR1)
% The averaging and angular projection cancel by their adjoint
% relation, and the entire reference variation cancels by this trace
% identity. Thus i Cdot=[h[C],C] is the native CAR Hartree evolution
% generated by E on this occupation/solution chart, and Edot=0 for
% a closed system. This constructs the derivative, not a claim that
% every proposed positive source has a unique global branch.
%
% In particular D=C-P(h) does NOT obey the old common-transport law:
%         Ddot=-i[h,D]-Pdot,   Pdot=DP_h[hdot].                 (RR2)
% The finite-radius mass variation retains this reference term;
% only the total boundary contact with h vanishes. Created pairs
% relative to this continued reference therefore remain in the
% source. No reference refilling or resetting of the physical C occurs.
%
% RR2. Exact radial response and its native pullback. In the dV measure
% let K be the symmetric geometric Hessian from section 7. Writing
% R=max(r,s) and I(R)=int_R^infinity 4pi G^2 J_t/(N_t u_t^3) dt,
% K_00(r,s)=-N_r N_s I(R),
% K_01(r,s)=G Theta(s-r)N_r/(s u_s)-u_s N_r N_s I(R),
% K_10(r,s)=K_01(s,r),
% K_11(r,s)=G Theta(r-s)N_s u_s/(r u_r)
%            +G Theta(s-r)N_r u_r/(s u_s)
%            -u_r u_s N_r N_s I(R).
% Equivalently its quadratic form is
% int dr[8pi Gr N delta M delta J/u
%                         -4pi G^2 N J(delta M)^2/u^3].
% The center limits are taken through the regular solution, and
% actual cell weights integrate the kernel, with Theta(0)=1/2.
% Put A=W^dag K W, so A is a real symmetric native matrix. This is
% the full density-current lattice kernel before reference response.
% The continued-reference susceptibility is
% Pi_IJ=2 Re sum_(a occupied,b empty)
%        <a|B_I|b><b|B_J|a>/(epsilon_a-epsilon_b).              (RR3)
% It is symmetric. It is negative semidefinite for the ground
% occupation, but need NOT have that sign after occupied levels have
% been continued through crossings. Only its symmetry is needed.
% With b_I=Tr(B_I delta C), delta q=b-Pi delta phi and
% delta phi=A delta q. On a regular chart with I+A Pi invertible,
%       K_c=(I+A Pi)^(-1)A=A(I+Pi A)^(-1)=K_c^T,
%       delta h=sum_I B_I(K_c b)_I.                            (RR4)
% Thus RR4 is the exact affine-covariance Hessian of E. It includes
% the source/reference feedback instead of discarding it as a higher
% order effect. On a pure-Slater orbit the ordinary second derivative
% additionally contains Tr(h delta^2 C), as usual for a curved chart.
%
% These are constructive finite-parent domains: a smooth regular
% positive M/r^3 at the center, positive M away from it, positive N,
% a nonsingular radial boundary linearization, separated opposite
% occupations and the stated implicit inverse give a local smooth
% E[C]. With analytic source maps the same finite-dimensional
% implicit construction is analytic. If these bounds are uniform,
% their limits can be taken; no uniform infinite-parent occupied
% separation, collapse interval or ell-independence is inferred.
%
% RR3. First orders in g=G/a^2, retaining mu=ma. All vertices and
% energies can be made dimensionless with a h0, a B_I and a E;
% the following expansion then uses g. At fixed normal sources the
% regular flow is order sqrt(g), so write the ACTUAL native response
% A=sqrt(g) A_(1/2)+g A_1+O(g^(3/2)), Pi=Pi_0+O(sqrt(g)).
% The complete response through first order g is
% K_c=sqrt(g) A_(1/2)
%         +g[A_1-A_(1/2) Pi_0 A_(1/2)]+O(g^(3/2)).             (RR5)
% A_(1/2),A_1 are obtained by the closed K and W above on the
% chosen background; they retain the density dependence of the flow.
% In a weak static particle scaling J=O(sqrt(g)), the lower-order
% density-current contribution is itself order g. In the controlled
% slow-particle and ell/geometric-length limit its density kernel
% reduces to -g mu^2/max(r,s). At finite ell the W factors and the
% full radial response remain: the once-filter Newton example in
% the older compact is not silently its exact finite-ell substitute.
%
% RR4. An actual ordered quadratic operator. Freeze a regular Hartree
% background Cbar, Pbar=P(hbar) and K_c. Define
% X_I=dGamma(B_I)-Tr(Pbar B_I), dbar_I=Tr[B_I(Cbar-Pbar)].
% With every reference contraction removed, set
% H_red^(2)=Ebar+sum_I hbar_I (X_I-dbar_I)
%       +(1/2)sum_IJ K_c,IJ : (X_I-dbar_I)(X_J-dbar_J) :_(Pbar).
% Here hbar=sum hbar_I B_I. The scalar and linear terms are the
% Taylor centering about the actual positive-mass background, not
% an omitted reference-sector energy. Equivalently the pair is
% (1/2)sum K_c,IJ :X_I X_J:_(Pbar), accompanied by the explicitly
% expanded linear counterterm -sum_IJ K_c,IJ dbar_J X_I and its
% scalar counterpart. This is a bounded self-adjoint operator on the
% ENTIRE finite-parent Fock space, with no imposed operator Gauss
% equation. Its use as a physical expansion is restricted to the
% stated positive Hartree chart and small actual fluctuations.
%
% For a Gaussian C, D=C-Pbar, the uncentered pair has exactly
% E_H^(2)=(1/2)sum K_c,IJ Tr(B_I D)Tr(B_J D),
% E_F^(2)=-(1/2)sum K_c,IJ Tr(B_I D B_J D),
% h_F^(2)=-sum K_c,IJ B_I D B_J.                              (RR6)
% Full centering changes only the scalar/linear Hartree terms.
% Ignoring the exchange correction, variation of the centered
% expectation gives hbar+sum B_I[K_c b(C-Cbar)]_I, the Taylor
% expansion of the EXACT Hartree generator RR1. The nonlinear
% Hartree construction itself uses E[C], not a frozen hbar.
% No band, spin or valley multiplicity is added to these full native
% matrices. For a single particle D=|psi><psi|, each Wick pair
% obeys Tr(B_I D B_J D)=Tr(B_I D)Tr(B_J D), so its self-attraction
% cancels exactly, including current and reference-response vertices.
% The uncentered interaction and its one-body derivative both vanish
% at D=0. This does not evaluate the positive-background Taylor
% approximation at the zero-mass configuration outside its chart.
%
% RR5. Physical mass variance, not a subtracted variance. Define the
% direct native mass row
% L_r q=N(r)^(-1) int_0^r dV_s N_s[(Wq)_rho+u_s(Wq)_J].
% Its full covariance response is
% ell_r=L_r(I+Pi A)^(-1),  A_r^eff=sum_I ell_(r,I) B_I.        (RR7)
% The linear fluctuation operator centered at Cbar has Gaussian
% physical variance V_r=Tr[Cbar A_r^eff(1-Cbar)A_r^eff].
% For Cbar=Pbar+Gamma with a Slater particle projector Gamma<=1-Pbar,
% V_r=V_(Pbar)+Tr[Gamma A_r^eff(1-2Pbar)A_r^eff]
%                           -Tr[Gamma A_r^eff Gamma A_r^eff].
% The excess is displayed beside the positive sea contribution; it
% need not be positive. In particular V_r<=V_(Pbar)+N_f||A_r^eff||^2.
% The controlled number approximation A=m F, 0<=F<=1 on the particle
% subspace, gives Var_particles<=m< M > and hence 1/sqrt(N_eff)
% relative fluctuations. It is not a bound on the preceding full
% sea-retaining variance for arbitrary relativistic states.
% Likewise in a positive rest-mass kernel with 0<k_min<=k<=k_max
% on the occupied region, |E_F|/|E_H|<=k_max/(k_min N_f) for a
% Slater projector. The exact RR6 is used when that domain fails;
% an indefinite density-current kernel has no universal such ratio.
%
% RR6. Domain and remainder retained explicitly. On a finite parent
% every fixed Taylor order, with the same complete reference Wick
% prescription, is a well-defined bounded Hermitian operator. A
% positive Hartree chart supplies its coefficients. Classical third
% derivatives give the usual O(||delta C||^3) Hartree remainder
% while the source remains in that chart. Its quantum validity also
% needs higher physical fluctuation moments: a second-moment bound
% alone is not a convergence theorem for the full ordered series.
% For an infinite parent or an all-orders Hamiltonian, a common
% invariant domain and convergence of the ordered higher terms must
% be supplied. RR1--RR7 do construct the Hartree functional derivative,
% its complete quadratic ordered correction and the actual variance;
% they do not manufacture those missing uniform/all-orders estimates.
% In particular no compression, no nonnormal-state premise and no
% point-field replacement is used to obtain the displayed operator.
%
% RR7. Accounting at finite radius. Along the actual Hartree motion,
% differentiating the solved mass functional gives
% Mdot_r=-i Tr(A_r^eff[h,C])=-i Tr(C[A_r^eff,h]).
% Equivalently use the direct mass vertex L_r B on
% Ddot=-i[h,D]-Pdot; the reference term is essential. At infinity,
% RR1 gives the gradient h and Edot=0 for a closed system. Open
% receiver/boundary fluxes therefore require this same charge
% allocation and an actual outgoing-state calculation; the former
% transported-reference work defect is not a binding back-reaction.
%
% SOURCES / INDEPENDENCE. Direct differentiation of the declared
% coherent radial rows, the exact adjoint pairing, P^2=P and [P,h]=0,
% spectral divided differences, CAR Wick contractions and finite
% Hermitian Taylor operators supply this derivation. The previously
% independently derived section-7 kernel and section-8 spectral
% formulas were rechecked and extended here to continued occupations
% and the actual full W^dag K W coupling. Current GOAL and both
% REVIEWA_REVIEW files were read completely, with owner overrides.
% No revised counterpart source was read before this substantive
% save. No new run, cross-lane X or full-series domain is asserted.

% RR8. TASK-8 BOOKKEEPING RE-DERIVED FOR THE SAME CONTINUED REFERENCE.
% Fix a finite parent, a sector reference P, a frozen one-body h and
% real symmetric pair coefficients V_IJ. Put p_I=Tr(P B_I),
% X_I=dGamma(B_I)-p_I, and
% H_sym=(1/4)sum_IJ V_IJ {X_I,X_J}.
% For Gaussian C=P+D, direct CAR contraction gives
% <H_sym>=E_H(D)+E_F(D)+F(P)+Tr(Sigma_P D), where
% F(P)=(1/2)sum V_IJ Tr[P B_I(1-P)B_J],
% Sigma_P=(1/2)sum V_IJ B_I(1-2P)B_J.
% The symmetric V sum makes the last matrix Hermitian. Thus the
% FULL reference-Wick pair is the exact operator identity
% H_W=H_sym-dGamma(Sigma_P)+c(P),
% c(P)=Tr(P Sigma_P)-F(P)
%     =-(1/2)sum V_IJ Tr(P B_I P B_J).                       (RR8)
% The identity follows either by normal ordering all four CAR
% operators or by the displayed independent constant/linear/quartic
% contractions. Normal ordering the one-body piece also subtracts
% Tr(hP). If a producer uses
% H_raw=dGamma(h)+H_sym-dGamma(Sigma_P),
% its missing scalar is
%             k(P)=c(P)-Tr(hP).                              (RR9)
% Comparing sectors with separately frozen P_N and P_0 therefore
% requires k(P_N)-k(P_0), using EACH sector's actual h and V if these
% also vary. This scalar leaves each fixed-sector eigenvector, energy
% gap and mode-cut entropy unchanged. It changes cross-sector total
% energies. A current producer must be compared to this exact
% operator convention; no claim about an unread new version follows.
%
% The independent variance correction is just as exact. For every
% scalar reference subtraction, Var[dGamma(A)-Tr(PA)]
% =Var[dGamma(A)]; in a Gaussian state it is Tr[C A(1-C)A].
% Var_C-Var_P is the signed excess, not that physical variance.
% sqrt(max(Var_C-Var_P,0)) both subtracts the sea and clips negative
% excess; it cannot be labelled the physical mass fluctuation.
% Full Wick subtraction in the interaction ENERGY removes reference
% contractions there, not the measured variance of this bilinear.
% These two corrections apply to the same reference/sector convention
% regardless of whether a fixed scalar was included in an eigensolver.
%
% ACTUAL RQ1--RQ4 COMPARISON / SOURCES.
% Read the complete saved original RED RQ1--RQ4 and SOURCES only after
% the substantive RR1--RR7 save and self-read. Its explicit W0/W1
% radial vertices give precisely W^dag K W here. Its scalar-clock
% integrating-factor variation and off-block spectral trace proof
% independently agree with RR1/RR2. The polarized radial quadratic
% form equals the four closed K blocks above; the superoperator
% (I+S DP)^(-1)S is RR4 in native source coordinates. Continued
% occupations need symmetry, not the ground-state sign of Pi.
% Its explicit c_H=Ebar-Tr(hbar Dbar)+Tr[Dbar L(Dbar)]/2 equals
% RR4's expanded centering constant; both reproduce the Hartree
% Taylor functional and the same Fock term. Its finite-parent
% implicit-function/commutator construction supplies local physical
% CAR evolution from an ACTUAL regular initial solution; no new
% initial solution or uniform parent/ell interval is assumed.
% First-g contact, rank-one cancellation, finite-polynomial domain
% and the open full-series positive-mass domain agree. The direct
% RR8--RR9 contraction was re-derived here for the lead's two task-8
% corrections. No new numerical result, external X or extra
% strong-operator continuation is used.

% ACTUAL RQ5 FIRST-g / BOOKKEEPING CHECK.
% Read the complete actual RQ5 after RR8--RR9 were saved. Independently
% substituting epsilon=sqrt(g) into the checked radial rows gives
% M0=int_0^r w rho, u1=-sqrt(2M0/r), M1=int_0^r w u1 J,
% u2=u1 M1/(2M0), n1=4pi int_r^infinity sJ/u1 ds.
% The source-mass vertices satisfy
% A0_r=int_0^r w B_rho,
% A1_r=int_0^r w[(n1(s)-n1(r))B_rho(s)+u1(s)B_J(s)]ds.
% With a_k(X)=Tr(Ak_r X), j(X)=Tr(B_J(r)X), the source-response
% coefficients in the RED's superoperator convention are
% Tr[Y S1(X)]=int dr{4pi r/u1 [a0(X)j(Y)+a0(Y)j(X)]
%                                  -4pi J a0(X)a0(Y)/u1^3},
% Tr[Y S2(X)]=int dr{4pi r/u1 [(n1-u2/u1)
%       *(a0(X)j(Y)+a0(Y)j(X))+a1(X)j(Y)+a1(Y)j(X)]
%    -4pi J/u1^3[(n1-3u2/u1)a0(X)a0(Y)
%                            +a1(X)a0(Y)+a0(X)a1(Y)]}.
% These follow by expanding N/u and N/u^3 in the polarized Hessian;
% the second term has no extra radial r. They agree term by term
% with RQ5 and give the explicit first-g response
% S2-S1 DP_flat S1. DP_flat is the spectral DERIVATIVE, not P0.
% At a source background varying with g its own Taylor terms must
% also be included. RQ5's independent scalar c(P), k(P) and physical
% variance decomposition agree exactly with RR8--RR9.

% ACTUAL CURRENT TASK-8 PRODUCER COMPARISON, NO NEW RUN.
% Read all of engine_hf_cluster_v4.py and engine_massfluct_v3.py at
% <internal>/engine, then the
% actual imported C1 chain/Hartree/sector routines, C2 target and
% comparison functions, engine_ops_v1.py, and the called fixed-
% reference/energy routines in engine_hf_cluster_scfcheck_v1.py and
% engine_hf_cluster_diag_v1.py. The following conclusions concern
% the code actually read, not an inferred run or numerical tolerance.
%
% In v4 exact_with_state, h1 is dGamma(h)-dGamma(Sigma_P), B_n
% is dGamma(hcal_n)-e_ref_n, and the matvec adds
% (1/2)sum_n B_n sum_m V_nm B_m. It contains neither c(P) nor
% -Tr(hP). run_case calls it once with H0["C0"] and separately
% with each H["C0"], then sets dE_ex=Ex-E0x. Therefore the exact
% energy differences still require RR9's k(P_N)-k(P_0). The
% docstring's claim that this c-number is sector-independent is
% not justified by the actual separately frozen references.
% The correction is to the scalar total-energy comparison; it
% changes neither those exact-sector states nor their mode-cut
% entropies. The actual pair is one two-component kappa/r chain
% with V_nm=-g/max(R_n,R_m), not the full RR4 native response.
%
% v4's ex_wick and sigma_wick implement the correct D-quadratic
% contractions. Its usual HF iteration recomputes the Hartree
% reference from its current potentials; the level-shift fallback
% instead fixes H["C0"] and uses DG.E_wick at that reference.
% The convergence/route labels do not make these two reference
% conventions identical. A variational comparison to one frozen
% exact operator uses its SAME reference, kernel and scalar origin.
%
% massfluct_v3 now separately computes the physical and reference
% variances for its ACTUAL unclocked cumulative energy vertex:
% V_C=F_mat+F_sea, V_P0=F_sea+B,
% delta V=F_mat-B, V_matter=F_mat+B,
% W_product=-Tr(P_m A P_m A).
% Direct CAR expansion verifies each identity, also for fractional
% positive-mode occupations. Its sqrt_V_C_over_E is the physical
% quantity; signed_sqrt_dV_over_E is only the signed excess and
% v2_fluct_mat_over_E is a different decomposition term.
% Thus version 3 has the required physical-variance column; a
% report must use that column with its sea/excess companions,
% rather than label the excess as physical. The one-particle
% same-kernel cancellation is present algebraically in its ExW
% and pair_energy functions, not claimed numerically evaluated here.
%
% This routine uses the chain's unclocked sharp cumulative energy,
% ground filling of the clocked chain, independent 16*kappa copies
% and an unsmoothed monopole kernel. It does not evaluate RR7's
% continued-reference/current/adjoint mass-response vertex. Those
% observable and native-3D limits remain explicit when using it as
% a control. No new output, current live job or numerical agreement
% has been inferred from a version name or a status row.

% REQUIRED NATIVE TASK-8 INTAKE: independent disposition before new RED
% comparison. Critical-first addendum and STATUS native_T8 mapping read;
% this serves F6/I2, paper 03_model.tex:382--385 and 05_star.tex:180--185.
% Read the complete actual solve/strand_E_spectral_task8_control.tex
% T8.1--T8.12, its review and delivered scope, including all SOURCES/endings,
% in <internal>/. Supporting reads: actual SB1--SB13
% source/review and SOURCES; FC4--FC5 canonical conventions and LH4--LH5
% exact two-native-hole convention, with their defining context/SOURCES.
% The older FC/LH reference dynamics is not imported into T8's continued
% spectral readout. No new counterpart comparison was read before this
% substantive save. This is an argument/specification check, not a run.
%
% At common coherent-sign/HA-adjoint/continued-reference premises, SB6 and
% T8.1 agree with RR1--RR4: L=(I+DP K)^(-1), K_c=K L=L^* K.
% T8's g denotes sqrt(G), not this note's dimensionless G/a^2; its expansion
% parameter is a sqrt(G/a^2). The same fixed-source radial Hessian, complete
% occupied/unoccupied Riesz response and actual signed Dbar are retained.
% T8.2's order-G coefficient includes the current-current Hessian, the
% displacement Dbar=nu F-g Sigma+O(g^2), and -K1 DP0 K1. Its positive
% prefactor on j_Sigma A tensor A/(r^2 b^3) follows from that negative
% displacement and the third derivative of the order-g boundary term;
% the signed b<0 is retained. This matches RR5 and the RQ5 comparison,
% including its background-shift qualification. The displayed JJ Gram
% form has no extra factor two and fixes no sign for the complete exchange.
%
% SB12's scalar and linear Taylor contacts equal RR4's c_H and
% hbar-mathcal L(Dbar), with mathcal L the covariance Hessian map;
% its full-reference Wick contraction equals RR6/RR8.
% T8.4--T8.6 retain the actual canonical cumulant
% -tau(F tensor F)(I-S), tau=nu(1-nu)/(D_+-1), rather than a Gaussian
% substitution. The signs of both initial covariance corrections and
% their source pullback L agree with the full CAR hierarchy. Pair
% self-cancellation applies to rank-one excess orthogonal to Rbar,
% not automatically to the signed Dbar. Absolute polynomial energies
% retain their scalar; RR9's k(P_N)-k(P_0) remains necessary if a raw
% centered-density implementation instead compares different references.
% No sector-dependent constant is licensed to cancel between sectors.
%
% T8.7 is RR7's current-inclusive mass selector with the full spectral
% pullback. At the outer boundary A_R=hbar since DP_h[hbar]=0.
% T8.8 is physical variance: its sea term is (1-nu)||P A_r F||_2^2
% and its canonical particle term is
% D_+ tau ||F A_r F-Tr(F A_r)F/D_+||_2^2. These coefficients follow
% from FC4--FC5 with the ordinary unnormalized antisymmetrizer.
% Fixed native number is not fixed Gauss energy; no sea variance is
% subtracted and no universal inverse-square-root number law is asserted.
%
% The SB complex contraction supplies an actual positive-mass/gapped
% finite-parent source chart and local Hartree solution for C_nu.
% The T8.9--T8.11 first-exit, second-time-derivative and Cauchy bounds
% compare exact polynomial, HF and nonlinear Hartree readouts on that
% same chart; frozen-polynomial energy is distinct from E(C_Q(t)).
% Their constants and zero-denominator conventions check against SB.
% On the specified open 5x1x1 parent, d=20 and N=18 make the complete
% physical block wedge^18 h exactly a 190-dimensional two-native-hole
% space. T8.12 retains the filled-state scalar, negative transposed HF
% matrix and pair coefficient one. Its initial two-hole density is
% Pi_(wedge^2 F^T)/45. This does not propagate or compress a Fock reference.
% The stated endpoint zero-mode and opposite-beta spectator tests retain
% the negative Dirac complement and a genuine pair coefficient on this
% thin parent. They are analytic nontriviality checks, not evaluated
% exact-versus-HF discrepancies.
%
% Disposition: no mandatory discrepancy found. This supplies matched
% other-route analytic support for the reduced gradient/Hessian, first-G
% contacts, full Wick/scalar and physical-variance components, plus a
% concrete native finite-matrix comparison specification. The counterpart
% retains its S/scoped-review status; no new numerical agreement is claimed.
% Its ma=1, ell=4a, R=8a parent is an implementation control, not a
% macroscopic star/core, continuum or averaging-independence calculation.
% The paper-facing one-page construction needs no repair from this intake.
% The required evolution/error checks remain the solve lane's stated
% control; the open all-orders domain is not made an added delivery gate.
% Published core unchanged; no new derivation, code, run or foreign copy.

% REQUIRED TASK-8 REFERENCE CHECK: independent disposition before new RED
% comparison. STATUS names this supplied control for F6/I2 and the first-
% correction sentences in 03_model.tex:382--385 / 05_star.tex:180--185.
% Actual complete hf_cluster_refcheck_v1.md and .json, complete producer
% solve/code/engine/engine_hf_cluster_refcheck_v1.py, and the called chain,
% Hartree, Wick, fixed-reference and level-shift functions were read in
% <internal>/. Necessary actual input rows/header
% of hf_cluster_v3_wick_scan2.json were matched. No new comparison from
% the original RED was read before this substantive disposition.
%
% The report supplies four selected auxiliary radial-chain comparisons:
% n_s=10, kappa=1, m=0.5, a=1, two components per site, R_n=n+1 and
% V_nm=-G/max(R_n,R_m). Three cases use G=0.12519591165303257 with
% N_f=4,7,8; the fourth uses G=0.03804863316614245 with N_f=8.
% This is the existing unsmoothed chain, not the native 3D source response
% or the separately supplied T8 190-dimensional native specification.
%
% The fixed solve freezes the separately rebuilt Hartree reference C0H.
% The moving solve recomputes its Hartree sea from the current potential.
% compare() evaluates BOTH determinants with DG.E_wick(C-C0H), so the
% displayed E_W comparison has one reference and one scalar origin within
% each case. The moving own-reference energy is separately reported.
% The reported projector Frobenius differences range from 3.98e-8 to
% 5.72e-6; the common-zero energy differences are at most 2.14e-11.
% The reference differences are nonzero (9.05e-6 to 2.31e-4 in Frobenius
% norm). Thus these data directly supply the previously absent comparison
% with an actual fixed-reference solve for these four cases, and show
% close projectors/common-zero trial energies there. They do not identify
% the two reference laws, functionals or their minima.
%
% All four pairs have reported convergence flags. The ordinary unshifted
% cases take 48--50 iterations and use the potential/Fock residual stopping
% rule followed by final diagonalization. The strong N_f=8 exception uses
% sigma=0.05: fixed and moving loops take 39 and 28 iterations respectively.
% The moving exception is seeded with the fixed solution. Crucially, the
% actual JSON records aufbau=false for BOTH exceptions, although their
% unshifted-Fock max-entry commutators are 2.69e-13 and 7.81e-13 and their
% reported Fermi gaps are about 0.00454. These are numerically commuting
% trial determinants; a positive Fock gap and small commutator do not
% certify lowest-orbital filling or a local/global variational minimum.
% No minimization certificate is supplied for the ordinary cases either.
% The inner moving-reference fixed-point routines have finite stopping
% tolerances but return no separate convergence flag in these records.
%
% The same-reference scalar cancels in a same-sector exact-minus-trial
% gap. This new comparison neither changes the raw exact energy assembly
% nor supplies RR9's cross-sector k(P_N)-k(P_0). Re-evaluating a moving
% determinant at C0H is a valid fixed-functional trial evaluation, not
% evidence that its moving-reference iteration minimizes that functional.
% This report evaluates no physical mass variance and changes none of
% the earlier sea-retaining versus excess-variance qualifications.
%
% Disposition: the absence of an explicit fixed-reference comparison is
% resolved at precisely these four reported chain cases. The exception's
% non-Aufbau diagnostic and all reference/minimization qualifications
% above remain necessary; no extension to the full scan is inferred.
% No mandatory paper-facing mathematical repair or compact change follows.
% The supplied report is a scoped numerical control, not a full-native
% validation, a cross-sector scalar repair, or a new all-orders-domain gate.
% No scientific code/run, foreign copy, new file or proof was made.

% REQUIRED STATIONARY COMPARISON, SC1--SC6 [R; independent new derivation,
% actual original-pair comparison below]. STATUS names 06_scaling.tex:84--89 before
% this work. SN1--SN14, its review and the stationary scaling handoff were
% ALREADY READ; their averaged-back target is known, not a blind prediction.
% No new original-RED stationary derivation was read before this save.
% The distinct route below uses the nonlinear invariant-subspace graph,
% direct coefficients of BOTH continuum Gauss rows and its orthonormal
% two-state effective matrix. It does not use SN's boundary Taylor/Hessian
% argument or source-scaling identity, a Newton half-counting argument, or
% a frozen quadratic replacement for the nonlinear boundary functional.

% SC1. Keep exactly SN's fixed a,m,ell>0, open cube |x_i|<=L with L/a even
% and >=2, full native h0=m beta+alpha_i p_i, Q_i and endpoint convention.
% Put Lambda=sum p_i^2, phi=phi_0, Gamma=|phi><phi| tensor 1_(beta=+1),
% Q=I-Gamma, and R=Q(m-h0)^(-1)Q. Gamma is the full isolated rank-two
% cluster at m, and phi is real, constant on the even sublattice and zero
% elsewhere. The normalized HA cell kernel, angular source projection,
% all source/field collars and coherent continuum Gauss rows are unchanged.
% Write dV=4pi r^2 dr, rho0=rho_Gamma, M0(r)=int_0^r rho0 dV and
% W(r)=-sqrt(2M0(r)/r), smoothly continued at r=0. M0(infinity)=2m.
% For the two field allocations use T_a=A_ell^* and T_p F=F|_Lambda;
% they act on scalar N and Cartesian components of v rhat. ONLY this
% field allocation changes in the p control: its source average is HA too.
% Set t=sqrt(G); the analytic extension uses signed t and
% u=-t sqrt(2M/r). Retain the actual source-dependent fields and define
% h_s(t,D)=e(T_s N_D)+P(T_s(v_D rhat)), s in {a,p}.
% On Gamma, J=0 exactly, N=1 and M=M0 at every t, so
% h_s(t,Gamma)=h0+t V_s, V_s=P(T_s(W rhat)). V_s=V_s,sp tensor I_4
% is Hermitian and purely imaginary in the spatial basis. All symbols
% here are the finite native matrices, with no extra valley factor.

% SC2. A second construction of the actual stationary branches. Represent
% a nearby rank-two projector by its graph Z:Ran Gamma -> Ran Q:
% F_Z=(Gamma+Z)(I+Z^dag Z)^(-1/2), D(Z)=F_Z F_Z^dag. The exact nonlinear
% stationarity equation, with h=h_s(t,D(Z)), is
% Qh Gamma+QhQ Z-Z Gamma h Gamma-Z Gamma hQ Z=0.
% At (t,Z)=(0,0) its real derivative in Z is (Qh0Q-m)Z, invertible by
% the isolated cluster gap; the derivative of h with respect to D is
% zero at t=0. The positive central HA density gives the common analytic
% Gauss branch on a neighbourhood of Gamma (the same BH/SN branch).
% Finite-dimensional analytic implicit inversion therefore supplies a
% unique small graph Z_s(t) for BOTH allocations. Shrink |t| so the
% Gauss branch, positive clock, gap at zero and selected cluster gap persist.
% D_s(t)=D(Z_s(t)) is then the actual rank-two spectral cluster of its
% OWN h_s(t,D_s(t)), not a state of a frozen Taylor matrix.
% Its leading graph and source are
% Z_s=t A_s+O(t^2), A_s=R V_s Gamma,
% D_s=Gamma+t D1_s+t^2 D2_s+O(t^3), D1_s=A_s+A_s^dag.
% Time reversal maps t to -t, fixes Gamma and reverses the native current.
% Uniqueness gives rho_s(t)=rho_s(-t), J_s(t)=-J_s(-t). Thus rho1_s=0,
% J_s=t j_s+O(t^3), where j_s(r)=Tr[B_J(r)D1_s] uses the SAME HA radial
% source vertex B_J in both controls. Scalar charge and level sum are
% analytic even functions of t. The estimates are on a fixed finite
% parent; no uniform bound in L/a, ell/a or m->0 is asserted.

% SC3. Direct nonlinear Gauss coefficients, with no boundary-gradient
% premise for point coupling. The exact rows are
% M'=4pi r^2(rho+uJ), (log N)'=-4pi G rJ/u, u=-sqrt(2GM/r), v=Nu.
% Along either branch they give
% M_s=M0+t^2 M2_s+O(t^4), N_s=1+t^2 n_s+O(t^4),
% v_s=t W+O(t^3),
% M2_s'=4pi r^2(rho2_s+W j_s),
% n_s(r)=int_r^infinity 4pi x j_s(x)/W(x) dx.
% These are regular at the centre: W=O(r), j_s=O(r). Sources have the
% fixed compact HA support; n_s vanishes outside it. Therefore
% h_s(t,D_s(t))=h0+t V_s+t^2 U_s+O(t^3), U_s=e(T_s n_s).
% The graph normalization, without determining any off-diagonal D2_s,
% gives Tr(h0 D2_s)=Tr[A_s^dag(h0-m)A_s]=-S_ss, where
% S_rs=Re Tr(Gamma V_r R V_s Gamma), r,s in {a,p}.
% HA normalization with complete collars gives int rho2_s dV=Tr(h0 D2_s).
% Its adjoint pairing gives int W j_s dV=Tr(V_a D1_s)=2S_as,
% INCLUDING when s=p: replacing V_a by V_p here would be incorrect.
% Consequently the ACTUAL boundary charge for either control is
% E_b,s=2m+G(-S_ss+2S_as)+O(G^2).                         (SC3)

% SC4. The occupied-level sum needs its own point calculation. Expanding
% the orthonormal graph matrix F_Z^dag h_s F_Z through t^2 gives
% m I_2+t Gamma V_s Gamma
%       +t^2[Gamma U_s Gamma+Gamma V_s R V_s Gamma]+O(t^3).
% Gamma V_s Gamma=0 since phi is real and V_s,sp is imaginary Hermitian.
% The two cross terms and the normalization/kinetic term combine into
% ONE resolvent term, not two. Thus
% E_e,s=Tr(h_s D_s)=2m+G(S_ss+L_s)+O(G^2),
% L_s=Tr(Gamma U_s)=2m sum_x |phi(x)|^2 (T_s n_s)_x.      (SC4)
% This is a closed first-order coefficient: j_s is the explicit finite
% matrix trace in SC2 and n_s the displayed one-dimensional integral.
% It contains no unknown second-order state or self-consistent field.
% Integrating n_s against the actual smooth rho0 gives
% L_a|_(j_s)=int rho0 n_s dV
%   =-int M0 n_s' dr=(1/2)int W j_s dV=S_as.
% The central boundary term vanishes and n_s(infinity)=0.
% Define the point lapse discrepancy on its OWN branch by
% Delta_p=2m sum_x |phi(x)|^2[n_p(|x|)-(A_ell^* n_p)_x].
% Then L_p=S_ap+Delta_p, whereas L_a=S_aa. Hence the matched results are
% E_b,a=2m+G S_aa+O(G^2), E_e,a=2m+2G S_aa+O(G^2),
% E_b,p=2m+G(-S_pp+2S_ap)+O(G^2),
% E_e,p=2m+G(S_pp+S_ap+Delta_p)+O(G^2).                  (SC4a)
% In particular E_e,p need not have twice E_b,p's correction.
% This derivation uses neither an adjoint identity for h_p nor SN10.

% SC5. Native closed forms and the actual coupling differences. Put
% f_s=V_s,sp phi, so <phi,f_s>=0. On that spatial complement,
% (m-h0)^(-1)=-(m+h0)Lambda^+. Summing the two beta-positive spin states
% gives S_rs=-b_rs with b_rs=4m Re<f_r,Lambda^+ f_s>.
% In particular b_a:=b_aa>0 and b_p:=b_pp>0: at the central site the
% same two-step stencil gives f_s(0)=phi(0)/(4ia) sum_i(T_s(W rhat)_i)_(2ae_i),
% which is nonzero. For a, HA's even coordinatewise decreasing kernel
% and full collars give strictly inward axial averages; for p the
% sampled inward component is already strictly negative. This is the
% previously checked actual native nonzero test, not a new potential.
% The averaged result is exactly the known SN8/SN11 target, now obtained
% by the graph/Gauss/effective-matrix route with the full Dirac complement.
% The point coefficients and differences are explicitly
% E_b,p=2m+G(b_p-2b_ap)+O(G^2),
% E_e,p=2m+G(-b_p-b_ap+Delta_p)+O(G^2),
% E_b,p-E_b,a=G b_delta+O(G^2),
% b_delta=4m<f_p-f_a,Lambda^+(f_p-f_a)>=b_p+b_a-2b_ap>=0,
% E_e,p-E_e,a=G(2b_a-b_p-b_ap+Delta_p)+O(G^2),
% E_e,p-2E_b,p+2m=G[3(b_ap-b_p)+Delta_p]+O(G^2).          (SC5)
% Nonnegativity is a statement about the FIRST-G boundary-difference
% coefficient; strict positivity of b_delta is not asserted. A zero
% coefficient supplies no sign for the higher-order difference. At
% arbitrary fixed ell the sign of the point binding coefficient and
% level-difference coefficient is not inferred from b_a>0.

% SC6. Reference, units and transfer domain. On each small-t branch set
% R_s(t)=1_(h_s(t)<0), C_s(t)=R_s(t)+D_s(t). The gap at zero remains open,
% so R_s is exactly the occupied sheet continued from the flat sea along
% this parameter path, with no crossing/refill ambiguity. C_s is a pure
% CAR covariance with exactly two particles and zero holes relative to
% this reference; both commute with their fixed branch Hamiltonian.
% The nonlinear calculation used D=C-R throughout. The RR spectral
% contact is not discarded: this parameterization constructs C after
% solving its exact excess source. For a the RR implicit inverse is
% nonsingular for small t (I+O(t) on the fixed gapped parent), and its
% actual E[C] has the native Hartree gradient. No such gradient identity
% is assigned to p; p is the requested different-coupling control.
% Neither parameter preparation is claimed to be a physical collapse.
% All O(G^2) remainders follow from analytic finite-parent branches and
% time reversal; their constants may depend on the fixed a,m,ell,L,
% cluster gap, kernel/collars and positive-source neighbourhood. In
% lattice units mu=ma, g=G/a^2, put beta_rs=a^3 b_rs and delta_p=a^3 Delta_p.
% Then a E_b,a=2mu-g beta_a+O(g^2); every other displayed energy formula
% uses a E, g, beta and delta_p in the same way. Thus the comparison is
% stated versus the specified Cavendish G/a^2, with no value chosen.
% Hydrodynamic a<<ell<<L_geom, uniform parent/ell control, spatial stress,
% independent geometric evolution, isolated-star and collapse transfers
% remain separate required tests, not inferred from this fixed-parent
% stationary matter/scalar-constraint result. No all-orders quantum
% domain premise or numerical control was added to this computation.
%
% SOURCES / PROVENANCE. Current GOAL_alt.md, critical-first addendum
% and pre-existing STATUS mapping; actual 06_scaling.tex:84--89; own
% independently saved RR1--RR7 and paired RQ1--RQ5; actual known
% solve/strand_B_stationary_native_branch.tex SN1--SN14, its review and
% delivered scope, plus strand_B_stationary_scaling_handoff.tex; actual
% strand_E_hydrodynamic_sources.tex HA1--HA8 and
% strand_E_boundary_hessian.tex BH1--BH5 only for their continuum source
% branch/adjoint definitions, not their old transported-reference law.
% This new graph/direct-Gauss/level derivation was saved before intake
% of the original RED's new stationary comparison. Rendered core unchanged;
% no new file, code, numerical run, foreign data copy or new state family.

% ACTUAL ST COMPARISON / SOURCES. After both substantive independent saves,
% read the complete original RED ST1--ST5, its SOURCES and exact ending,
% and the current unchanged rendered quantum core. Its nonlinear graph
% branch, signed-t parity, direct mass/clock coefficients, full-complement
% spin trace and point lapse discrepancy agree with SC1--SC6. ST4.4 is
% the equivalent Fubini formula for Delta_p; its evaluation measure is
% not a new source, and its central atom has the stated integrable weight.
% No mandatory correction was found. Both original workers independently
% selected the graph/direct-Gauss route: this is ONE distinct route from
% the already known SN Riesz/boundary-Taylor/source-scaling proof, not two
% additional methods. The matched SN8/SN11 agreement is only the stated
% fixed-parent, massive, gapped averaged-back branch. Point coefficients
% are the new paired analytic control, with no numerical or hydrodynamic
% transfer claimed. Prior validated core/proofs remain unchanged; complete
% new argument/SOURCES/ending coverage and exact current-core read are done.
